\documentclass[a4paper]{article}
\usepackage[left=10truemm,right=10truemm,top=25truemm,bottom=20truemm]{geometry}
\usepackage{mathtools}
\usepackage{amsmath}
\usepackage{amsfonts}
\usepackage{bm}
\usepackage{setspace}
\usepackage{wrapfig}
\usepackage[dvipdfmx]{hyperref}
\usepackage{pxjahyper}
\usepackage{docmute}
\usepackage{amssymb}
\DeclareMathOperator\arctanh{arctanh}
\DeclareMathOperator\arccosh{arccosh}

\title{Double Spacetime Theory:\\ Gravity as Torsion between Usual and Dual Spacetime}

\author{https://github.com/hypernumbernet}
\date{\today}

\begin{document}

\maketitle

\begin{abstract}
General Relativity (GR) describes gravity as the curvature of a single spacetime continuum. Despite its empirical success, the theory encounters profound difficulties when confronted with quantum mechanics and the nature of singularities. This paper proposes a radical reformulation: gravity is not curvature of a continuous manifold, but torsion arising from the relative misalignment between two compact spacetimes intrinsically attached to every massive particle.

Employing the 16-real-dimensional biquaternion algebra isomorphic to Cl(3,1), we define the usual spacetime with basis (j, kI, kJ, kK) and the dual spacetime with basis (k, jI, jJ, jK). The dual map $X \mapsto Xi$ reverses the arrow of time while preserving the Minkowski norm $X^2 = X'^2$. Lorentz transformations and local frame rotations are generated by commuting pairs of rotors acting separately on each spacetime.

The complete rotor is $R_\text{total} = \exp[(\phi_1/2)iI + (\phi_2/2)iJ + (\phi_3/2)iK + (\theta_1/2)I + (\theta_2/2)J + (\theta_3/2)K]$.
The torsional mismatch rotor $\Omega = R^\dagger_\text{usual} R_\text{dual}$ yields the torsion bivector $\Omega_\text{biv} = \log \Omega$.
The unique Lorentz-invariant scalar is the Killing form on $\mathfrak{so}(3,1) \oplus \mathfrak{so}(3,1)$, defined as $J = \frac{1}{16} B(\Omega_\text{biv}, \Omega_\text{biv})$.

The action $S = \frac{c^4}{16\pi G} \int J \, d^4x$ is shown to be dynamically equivalent to the Einstein-Hilbert action, reproducing GR exactly while eliminating Christoffel symbols, Riemann curvature, and the continuum spacetime hypothesis altogether. The theory is the biquaternionic realization of the Teleparallel Equivalent of GR (TEGR), with torsion now physically interpreted as dual-spacetime torsion.

Crucially, since $J$ depends only on the relative angle between particle-intrinsic dual rotors, coherent excitation of the dual-spacetime components $\theta_a$ offers a direct pathway to gravitational engineering: gravity shielding, inertial mass reduction, and even antigravity become, for the first time, questions of rotor synchronization rather than violations of fundamental law.

GR is not superseded—it is completed, and gravity is rendered controllable in principle.
\end{abstract}

\section{Introduction}
\label{sec:introduction}

The theory of GR, since its inception in 1915, has stood as one of the most successful descriptions of gravitation in physics. Yet, despite its empirical triumphs—from the perihelion precession of Mercury to gravitational waves and black hole imaging—a profound philosophical unease persists: why does matter curve spacetime, and why does curved spacetime, in turn, dictate the motion of matter? The Einstein field equations tell us how gravity works with exquisite precision, but they offer no deeper reason for the origin of curvature itself. This explanatory gap becomes particularly acute when confronting Mach's principle: the inertia of a body should be determined by the totality of matter in the universe, yet GR treats spacetime as an independent, continuous entity that can exist even in complete emptiness.

A century of attempts to resolve these issues—loop quantum gravity, string theory, emergent gravity—has yielded mathematical elegance but no consensus on the microscopic origin of spacetime or the physical mechanism that makes mass-energy generate curvature. All these approaches retain the continuum hypothesis: spacetime is a smooth, differentiable manifold a priori, and gravity is an effectuated through the metric's deviation from flatness via Christoffel symbols and the Riemann tensor.

The present work breaks radically with this tradition.

We propose that the spacetime continuum hypothesis is fundamentally incorrect. There is no single, shared spacetime manifold. Instead, each individual particle carries its own pair of intrinsically double spacetimes—a compactified 16-real-dimensional structure encoded in the biquaternion algebra isomorphic to the Clifford algebra Cl(3,1). The usual spacetime (with basis j, kI, kJ, kK) and its dual (with basis k, jI, jJ, jK) are related by right-multiplication by $i$, a transformation that reverses the arrow of time while preserving the Minkowski norm. Gravity and inertia arise not from curvature of a background manifold, but from the torsional mismatch angle between these two particle-intrinsic spacetimes when parallel transport is demanded across different particles.

This philosophical shift is profound. In standard GR, spacetime is an autonomous entity that tells matter how to move, while matter tells spacetime how to curve (Wheeler's famous phrase). In the double spacetime theory, there is no autonomous spacetime at all. Spacetime degrees of freedom are demoted to internal attributes of particles themselves—much like spin or charge—and the apparent curvature of the macroscopic world emerges only as a collective effect of the torsional misalignment among the dual rotors of neighboring particles. The Einstein equations are recovered exactly, yet Christoffel symbols, covariant derivatives, and the Riemann tensor are entirely eliminated: they are replaced by simple algebraic operations (rotor exponentiation, logarithm, and the Killing form trace) on biquaternions.

The theory therefore achieves three revolutionary outcomes simultaneously:

1. Exact equivalence with GR in all observable predictions, including the Newtonian limit, black holes, cosmology, and gravitational waves—yet without any geometric curvature in the traditional sense.

2. Natural unification of gravity and inertia: the equivalence principle emerges as the requirement that the usual and dual rotors remain parallel in free fall; acceleration violates this parallelism, producing inertial forces as the same torsional phenomenon as gravity.

3. In-principle controllability of gravity: because the scalar invariant $J$ depends only on the relative angle between particle-intrinsic dual rotors, external fields capable of coherently exciting the dual-spacetime components $\theta_a$ can reduce, cancel, or even reverse gravitational interaction. Gravity ceases to be an inviolable fundamental force and becomes an engineerable degree of freedom.

This is not a modification of GR; it is its completion. The continuum was a useful fiction—beautiful, but ultimately unnecessary. Once spacetime is recognized as a paired, particle-local structure, the century-old mystery of why matter curves spacetime dissolves: matter does not curve spacetime; it misaligns its own double spacetimes, and this misalignment is what we perceive as gravity.

The remainder of this paper is organized as follows.

Section~\ref{sec:math} establishes the mathematical foundation by introducing the biquaternion algebra isomorphic to Cl(3,1) and constructing the dual spacetime pair intrinsic to every particle.

Section~\ref{sec:kinematics} develops the complete rotor formalism that unifies Lorentz boosts, rotations, and parallel transport without matrices or Christoffel symbols.

Section~\ref{sec:gravity} demonstrates that the torsional mismatch scalar $J$, constructed solely from the relative rotor $\Omega = R_\text{usual}^\dagger R_{\text{dual}}$, yields an action dynamically equivalent to the Einstein-Hilbert action, revealing the theory as the biquaternionic realization of Teleparallel Gravity.

Section~\ref{sec:unified} shows that inertia and gravity are identical phenomena arising from the same dual-rotor rigidity, with the equivalence principle emerging as an algebraic identity.

Section~\ref{sec:darkmatter} eliminates the need for dark matter by deriving flat galactic rotation curves from baryons alone.

Section~\ref{sec:strongfield} reinterprets stellar remnants as torsion stars with multi-layered attractive-repulsive structure, forbidding both singularities and event horizons.

Section~\ref{sec:timedilation} derives time dilation effects from torsional mismatch, predicting measurable deviations from GR in strong gravitational fields and planetary interiors.

Section~\ref{sec:nuclear} proves scale invariance of the torsional layers, identifying the nuclear force as gravity with layers, and atomic nuclei as primordial torsion stars.

Section~\ref{sec:electronshells} models the electron shell as an electrical inversion layer, explaining electron orbitals and interatomic forces from the perspective of microscopic-scale electromagnetism.

Section~\ref{sec:unification} outlines the unification of all fundamental forces within the biquaternion algebra, with gauge fields emerging as additional rotor components.

Section~\ref{sec:gravitycontrol} details practical mechanisms for gravitational engineering via coherent excitation of dual-spacetime rotors using circularly polarized electromagnetic fields.

Section~\ref{sec:baryonasymmetry} proposes a mechanism for baryon asymmetry based on dual-spacetime torsion dynamics in the early universe.

Section~\ref{sec:discretetime} introduces a discrete model of dual spacetime, where the torsional layers are quantized and bounded, leading to a finite unit group in the corresponding integer ring and enabling rigorous Diophantine constraints.

Section~\ref{sec:dualrotordynamics} develops the dynamics of dual rotors under external fields, deriving equations of motion and stability criteria for torsional configurations.

Section~\ref{sec:quantumgravity} extends the theory to quantum gravity, identifying quantum states with dual-spacetime rotor phases and deriving uncertainty relations from torsional coherence conditions.

Section~\ref{sec:conclusions} summarizes the unification of all forces within the 16-dimensional biquaternion algebra and outlines near-term experimental tests that can confirm or falsify the theory within this decade.

With this reformulation, GR is liberated from the prison of the continuum and elevated to a theory in which gravity is, for the first time, controllable in principle.

\section{Mathematical Foundation: The 16-Dimensional Biquaternion Algebra}
\label{sec:math}

The entire double spacetime theory is constructed within the 16-real-dimensional algebra of biquaternions (also known as double quaternions), which is canonically isomorphic to the Clifford geometric algebra $\mathrm{Cl}(3,1)$ over Minkowski spacetime with signature $(-,+,+,+)$. This algebra naturally accommodates two complete, mutually commuting copies of Minkowski spacetime — the usual and the dual — within a single algebraic structure intrinsically attached to every massive particle.

\subsection{The Biquaternion Algebra}

We begin with two independent copies of Hamilton's quaternion algebra:
\begin{itemize}
  \item Primary quaternions generated by $i,j,k$ obeying $i^2=j^2=k^2=-1$, \ $ij=k$, \ $ji=-k$ (cyclic),
  \item Secondary quaternions generated by $I,J,K$ obeying identical rules $I^2=J^2=K^2=-1$, \ $IJ=K$, \ $JI=-K$ (cyclic).
\end{itemize}
The full biquaternion algebra is their tensor product $\mathbb{H} \otimes_{\mathbb{R}} \mathbb{H}$, in which the two sets strictly commute:
\[
iI = Ii,\; iJ = Ji,\; iK = Ki,\; jI = Ij,\; \dots
\]
The resulting 16 linearly independent real basis elements are
\[
1,\; i,\; j,\; k,\; I,\; J,\; K,\; iI,\; iJ,\; iK,\; jI,\; jJ,\; jK,\; kI,\; kJ,\; kK.
\]

\subsection{Explicit Isomorphism with Cl(3,1)}

The algebra is isomorphic to $\mathrm{Cl}(3,1)$ via the following identification of grade-1 (vector) generators:
\begin{align*}
  e_0 &= j,                    & &\text{(time-like basis vector)}\\[4pt]
  e_1 &= kI, \quad e_2 = kJ, \quad e_3 = kK & &\text{(space-like basis vectors)},
\end{align*}
which immediately satisfy the defining relations
\[
e_0^2 = -1, \quad e_1^2 = e_2^2 = e_3^2 = +1, \quad \{e_\mu,e_\nu\} = 0 \;\; (\mu \neq \nu).
\]
The higher-grade elements then follow standard Clifford multiplication. In particular:
\begin{alignat*}{3}
  &\text{grade-2 (bivectors):}  &\quad& e_0e_1 = iI, \;\; e_0e_2 = iJ, \;\; e_0e_3 = iK, \\
  &&& e_3e_2 = I,   \;\; e_1e_3 = J,   \;\; e_2e_1 = K, \\[6pt]
  &\text{grade-3 (trivectors):} && e_1e_2e_3 = k, \;\; e_0e_3e_2 = jI, \;\; e_0e_1e_3 = jJ, \;\; e_0e_2e_1 = jK, \\[6pt]
  &\text{pseudoscalar:}        && e_0e_1e_2e_3 = i.
\end{alignat*}
The complementary ``dual'' basis (time-reversed under right-multiplication by $i$) is generated by
\[
\tilde{e}_0 = k = e_1e_2e_3, \quad \tilde{e}_1 = jI = e_0e_3e_2, \quad \tilde{e}_2 = jJ = e_0e_1e_3, \quad \tilde{e}_3 = jK = e_0e_2e_1,
\]
yielding the dual spacetime representation introduced below.

This explicit isomorphism makes manifest that the 16-dimensional biquaternion algebra is not an ad-hoc construction but the geometrically natural Clifford algebra of 4-dimensional Minkowski spacetime, now enriched with an intrinsic duality that distinguishes usual and dual sectors while preserving the full Lorentzian structure. 

The apparent irregularity in the Clifford algebraic forms — such as the asymmetric placement of $j$ as the sole time-like vector generator and the cross-term-dominated spatial basis $kI$, $kJ$, $kK$ — is not accidental: it exists precisely to preserve the full rotational symmetry of three-dimensional physical space within a structure that simultaneously encodes two time-reversed copies of Minkowski spacetime. This subtle asymmetry is what allows the dual map $X \mapsto Xi$ to reverse the arrow of time without breaking spatial isotropy, thereby providing the algebraic foundation for both the equivalence principle and the emergence of gravitational torsion as a relative misalignment between particle-intrinsic dual frames.

\subsection{Dual Spacetime Structure Intrinsic to Particles}

Each particle is postulated to carry an intrinsic pair of compactified Minkowski spacetimes encoded within the biquaternion algebra:
\begin{itemize}
  \item \textbf{Usual Spacetime:} Represented by vectors of the form
  \[
  X = ct \, j + x \, kI + y \, kJ + z \, kK,
  \]
  with Minkowski norm
  \[
  X^2 = -(ct)^2 + x^2 + y^2 + z^2.
  \]

  \item \textbf{Dual Spacetime:} Represented by vectors of the form
  \[
  X' = ct' \, k + x' \, jI + y' \, jJ + z' \, jK,
  \]
  with identical Minkowski norm
  \[
  X'^2 = -(ct')^2 + x'^2 + y'^2 + z'^2.
  \]
\end{itemize}
The two spacetimes are related by the dual map
\[
X' = X i,
\]
which reverses the arrow of time:
\[
(ct \, j) i = ct \, (j i) = -ct \, k,
\]
while preserving the Minkowski norm. Since the pseudoscalar $i = e_0 e_1 e_2 e_3$ of $\mathrm{Cl}(3,1)$ anticommutes with every grade-1 vector ($i X = -X i$ for $X \in \mathrm{Cl}(3,1)$) and satisfies $i^2 = -1$, the norm transforms as
\[
X'^2 = (X i)(X i) = X (i X) i = -X (X i) i = -X^2\, i^2 = -X^2 \cdot (-1) = X^2.
\]
This dual structure is intrinsic to every particle, with the usual spacetime governing its standard kinematics and the dual spacetime encoding complementary degrees of freedom that become dynamically relevant in the presence of gravity and inertia. The existence of these two spacetimes allows for a richer geometric interpretation of Lorentz transformations, parallel transport, and ultimately the nature of gravitation itself.

\section{Kinematic Structure: Rotors, Boosts, and Parallel Transport}
\label{sec:kinematics}

In this section we develop the kinematic framework of the theory. All Lorentz transformations and local frame rotations are generated by commuting pairs of rotors acting independently on the usual and dual spacetimes. The key insight is that boosts (traditionally hyperbolic rotations) arise naturally from bivectors that square to +1, yielding the familiar $\cosh$/$\sinh$ expansion of the rotor without any ad-hoc introduction of imaginary angles.

\subsection{Rotors in the Usual Spacetime}

To illustrate the kinematic structure, consider a Lorentz boost along the $x$-direction with rapidity $\phi$. The generating bivector is $iI$, which satisfies $(iI)^2 = i^2 I^2 = (-1)(-1) = +1$. In general, all bivectors in the usual spacetime ($iI$, $iJ$, $iK$) square to $+1$, imparting a boost-like character, while those in the dual spacetime ($I$, $J$, $K$) square to $-1$, evoking pure rotations.

The corresponding boost rotor is thus
\[
R_\text{usual} = \exp\!\left( \frac{\phi}{2} iI \right).
\]
Given the positive square of the generator, the exponential expands naturally into hyperbolic functions:
\[
R_\text{usual} = \cosh \frac{\phi}{2} + iI \sinh \frac{\phi}{2}.
\]
Applying this rotor via the sandwich product to a spacetime vector $X = ct \, j + x \, kI + \cdots$ yields the canonical Lorentz transformation:
\begin{align*}
ct' \, j + x' \, kI &= R_\text{usual} \, X \, R_\text{usual}^\dagger \\
&= \left[ \cosh\frac{\phi}{2} + iI \sinh\frac{\phi}{2} \right] (ct \, j + x \, kI) \left[ \cosh\frac{\phi}{2} - iI \sinh\frac{\phi}{2} \right] \\
&= (\gamma ct + \gamma v x) \, j + (\gamma x + \gamma v ct) \, kI,
\end{align*}
where $\gamma = \cosh \phi$ and $v/c = \tanh \phi$. Pure rotations within the usual spacetime arise from bivectors like $I$, $J$, $K$, which square to $-1$ and produce the familiar trigonometric expansions of $\cos$ and $\sin$.

\subsection{The Complete Rotor}

The most general proper orthochronous transformation on the full biquaternion structure is generated by the complete rotor
\[
R_\text{total} = \exp\left( \frac{\phi_1}{2} iI + \frac{\phi_2}{2} iJ + \frac{\phi_3}{2} iK + \frac{\theta_1}{2} I + \frac{\theta_2}{2} J + \frac{\theta_3}{2} K \right),
\]
where the parameters $\phi_a$ govern boosts and rotations in the usual spacetime, and $\theta_a$ those in the dual. Since the usual and dual generators commute (e.g., $[iI, I] = 0$), the rotor factorizes neatly as
\[
R_\text{total} = R_\text{usual} \, R_\text{dual}.
\]

To make the geometric meaning more transparent, we now introduce the bivector form of the rotor, which separates magnitudes from directions and eliminates explicit summation indices.

Define the unit boost bivector in the usual spacetime
\[
\hat{p} = p_1 \, iI + p_2 \, iJ + p_3 \, iK, \qquad \hat{p}^2 = +1,
\]
where $\hat{p}$ is normalized ($|\hat{p}| = 1$) and $\phi = \sqrt{\phi_1^2 + \phi_2^2 + \phi_3^2}$ is the total rapidity (with direction encoded in the components $p_a = \phi_a / \phi$).

Similarly, define the unit rotation bivector in the dual spacetime
\[
\hat{q} = q_1 \, I + q_2 \, J + q_3 \, K, \qquad \hat{q}^2 = -1,
\]
where $\hat{q}$ is normalized ($|\hat{q}| = 1$) and $\theta = \sqrt{\theta_1^2 + \theta_2^2 + \theta_3^2}$ is the total dual rotation angle (with direction encoded in $q_a = \theta_a / \theta$).

The complete rotor then takes the elegant bivector-exponential form
\[
R_\text{total} = \exp\!\left( \frac{\phi}{2} \hat{p} + \frac{\theta}{2} \hat{q} \right)
= R_\text{usual} \, R_\text{dual},
\]
with
\[
R_\text{usual} = \exp\!\left( \frac{\phi}{2} \hat{p} \right)
= \cosh\frac{\phi}{2} + \hat{p} \sinh\frac{\phi}{2},
\]
\[
R_\text{dual} = \exp\!\left( \frac{\theta}{2} \hat{q} \right)
= \cos\frac{\theta}{2} + \hat{q} \sin\frac{\theta}{2}.
\]

This unified bivector form encodes the full Lorentz group action on both spacetimes simultaneously while making manifest the six independent real parameters: the rapidity $\phi$ and unit boost direction $\hat{p}$ in the usual spacetime, and the rotation angle $\theta$ and unit rotation axis $\hat{q}$ in the dual spacetime.

\subsection{The Sandwich Product}

The versor formalism provides a compact means to apply these transformations: for any biquaternion multivector $X$ representing a spacetime event, the transformed vector is
\[
\tilde{X} = R_\text{total} \, X \, R_\text{total}^\dagger.
\]
This single operation preserves the Minkowski norm ($X^2 = \tilde{X}^2$) and acts coherently on both usual and dual components, seamlessly integrating boosts and rotations into one algebraic expression. No matrices or coordinate-dependent machinery are required; the geometry unfolds intrinsically within the algebra.

\section{Gravity as Torsional Mismatch between Dual Spacetime}
\label{sec:gravity}

In this central section, we establish that gravity is neither spacetime curvature nor an external force, but the torsional mismatch between the usual and dual spacetimes intrinsically carried by every massive particle. The Einstein-Hilbert action emerges algebraically from this mismatch, without Christoffel symbols, Riemann tensors, or any reference to a continuous manifold. The resulting theory is mathematically equivalent to the TEGR, yet provides a profound physical reinterpretation: torsion is now the relative rotation angle between particle-local dual rotors.

\subsection{Definition of the Torsional Mismatch Rotor}

The torsional mismatch rotor is defined as
\[
\Omega = R_\text{usual}^\dagger R_\text{dual}.
\]

The dagger operation is the standard conjugate-inverse of the rotor in the biquaternion algebra:
\[
R_\text{usual}^\dagger = \exp\!\left( -\frac{\phi}{2} \hat{p} \right)
= \cosh\frac{\phi}{2} - \hat{p} \sinh\frac{\phi}{2}.
\]
Although formally a simple conjugation, this dagger carries special significance in the double spacetime framework: the negative sign in front of the sinh term precisely encodes the time-reversal inherent in the dual map. Since the pseudoscalar $i$ is grade-4 and therefore commutes with bivectors, right-multiplication by $i$ acts on the boost generator as
\[
\hat{p}\, i = - (p_1 I + p_2 J + p_3 K) \equiv -\hat{q}',
\qquad (\hat{p}\, i)^2 = -1,
\]
where $\hat{q}' \equiv p_1 I + p_2 J + p_3 K$ is a unit rotation generator in the dual sector. Equivalently, $\hat{p} = \hat{q}'\, i$, and the inverse usual rotor admits the exact rewriting
\[
R_\text{usual}^\dagger = \cosh\frac{\phi}{2} - \hat{q}'\, i\, \sinh\frac{\phi}{2}.
\]
Using the standard identities $\cosh(\phi/2) = \cos(i\phi/2)$ and $\sinh(\phi/2) = \sin(i\phi/2)/i$, this admits the analytically continued ($\phi \to i\phi$) trigonometric form
\[
R_\text{usual}^\dagger = \cos\!\left(\frac{i\phi}{2}\right) - \hat{q}'\, \sin\!\left(\frac{i\phi}{2}\right),
\]
revealing that the dagger operation effectively bridges the hyperbolic geometry of the usual spacetime to the elliptic geometry of the dual spacetime---not as an equality of $\cosh$ and $\cos$ for real arguments, but as the unique analytic continuation of the boost rotor across the duality $X \mapsto X i$.

This definition is structurally reminiscent of operations in Inter-Universal Teichmüller theory: the dagger temporarily ``deconstructs'' the hyperbolic boost structure of the usual spacetime (inducing the dual map through its sign reversal), while subsequent multiplication by $R_\text{dual}$ ``reconstructs'' the combined object, making the relative misalignment manifest. Through this ``deconstruct-then-reconstruct'' process, the torsional mismatch between the usual and dual spacetimes is extracted algebraically, and this mismatch constitutes the essence of gravity in the double spacetime framework.

In the vacuum limit (free fall), perfect anti-synchronization requires $R_\text{dual} = R_\text{usual}^\dagger$ (i.e., $\theta = \phi$ and $\hat{q} = -\hat{q}'$), yielding $\Omega = 1$ and vanishing torsion. Mass-energy induces deviations from this synchronization, producing a non-trivial $\Omega \in \mathrm{Spin}^+(3,1) \oplus \mathrm{Spin}^+(3,1)$.

The associated torsion bivector is
\[
\Omega_\text{biv} = \log \Omega \in \mathfrak{so}(3,1) \oplus \mathfrak{so}(3,1),
\]
where $\log$ is taken as the principal branch defined on a simply connected neighbourhood of the identity $\Omega = 1$ in $\mathrm{Spin}^+(3,1) \oplus \mathrm{Spin}^+(3,1)$, the natural domain on which the Baker--Campbell--Hausdorff (BCH) expansion converges. Throughout this paper the regime of physical interest is the weak-field/small-mismatch sector $\Omega \approx 1$, where $\Omega_\text{biv}$ is uniquely determined; topologically distinct branches lying outside this neighbourhood are not used in the dynamics presented here.

\subsection{Lorentz-Invariant Scalar from Dual Torsion}

The Killing form on $\mathfrak{so}(3,1)$ (vector representation) is
\[
B(X,Y) = 4\,\mathrm{Tr}(X Y).
\]
Denote the three grade-2 spatial bivectors of \S\ref{sec:kinematics} collectively by $\Gamma_a$ ($a=1,2,3$) with $\Gamma_1 = I$, $\Gamma_2 = J$, $\Gamma_3 = K$ (the Clifford identification $e_3e_2 = I$, $e_1e_3 = J$, $e_2e_1 = K$ of \S\ref{sec:math}). The usual-spacetime boost generators are $i\Gamma_a$, while $\Gamma_a$ generate rotations in the dual sector. Evaluating the Killing form gives $B(i\Gamma_a, i\Gamma_a) = +8$ and $B(\Gamma_a, \Gamma_a) = -8$, reproducing the Lorentzian signature perfectly.

The unique quadratic invariant is
\[
J = \frac{1}{16} B(\Omega_\text{biv}, \Omega_\text{biv}) = \frac{1}{2} \mathrm{Tr}(\Omega_\text{biv}^2).
\]

$J = 0$ in vacuum, $J > 0$ for ordinary attractive gravity (usual-spacetime boosts dominate), and $J < 0$ is possible when dual-spacetime rotations dominate (repulsive gravity).

\subsection{Exact Equivalence to the Einstein-Hilbert Action}

We propose the action
\[
S = \frac{c^4}{16\pi G} \int J \, d^4x.
\]
This action is purely algebraic and background-independent.

The scalar $J$ is identical (up to normalization and boundary terms) to the torsion scalar $T$ of TEGR. It is rigorously proven in the literature that
\[
\int T \, e \, d^4x = -\int R \, \sqrt{-g} \, d^4x + \mathrm{boundary},
\]
so the vacuum and matter-coupled field equations are exactly the Einstein equations. All solutions of GR (Schwarzschild, Kerr, FLRW, gravitational waves, etc.) are recovered precisely.

Thus, GR is reproduced in full, yet gravity is demoted to a torsional mismatch angle between particle-intrinsic dual spacetimes.

\subsection{Weak-Field and Newtonian Limit}

In the linearized regime,
\[
R_\text{usual} \approx 1 + \frac{1}{2} \sum \alpha_a \, i\Gamma_a, \quad R_\text{dual} \approx 1,
\]
yielding $\Omega_\text{biv} \approx \sum \alpha_a \, i\Gamma_a$ and $J \approx \frac{1}{2} \sum \alpha_a^2$. Standard TEGR linearization shows $\alpha_a \propto \partial_a (\Phi/c^2)$, leading to Poisson's equation
\[
\nabla^2 \Phi = 4\pi G \rho
\]
with the correct coefficient. Post-Newtonian parameters are identical to GR ($\gamma = \beta = 1$, etc.).

\subsection{Rotor Action in the Dual Spacetime}
\label{subsec:dual-rotors}

In the double spacetime theory, the sole physical realities are the rotor angles within the intrinsic biquaternion algebra $\mathbb{H} \otimes_{\mathbb{R}} \mathbb{H} \cong \mathrm{Cl}(3,1)$ carried by every massive particle. Unlike string theory, which posits extended physical objects as fundamental entities, here the boost and rotation angles encoded in the commuting pairs of rotors are the primary ontological elements. Gravity, inertia, and quantum states all emerge from deviations and coherences among these angles.

The dual spacetime, related to the usual spacetime by the dual map $X \mapsto Xi$, inherits a time-reversed orientation. Consequently, an angle that generates a pure spatial rotation in the dual sector must be appropriately dual-transformed when incorporated into the unified rotor description.

Consider a dual-sector angle $\theta$. In the pure dual basis, it produces a rotation generated by the elliptic versor $\hat{q}$ ($\hat{q}^2 = -1$):
\[
\exp\left( \frac{\theta}{2} \hat{q} \right).
\]
However, since the dual map reverses the arrow of time, the physical effect of this angle on the dual spacetime is equivalent to a boost when viewed through the dual transformation.

To reflect this correctly in the torsional mismatch calculation, we apply the dual conversion to the angle itself: $\theta \mapsto \theta i$. This transforms the dual contribution into a hyperbolic boost term:
\[
\exp\left( \frac{\theta i}{2} \hat{q} \right) = \exp\left( \frac{\theta}{2} \hat{q}i \right) = \exp\left( \frac{\theta}{2} (q_1 iI + q_2 iJ + q_3 iK) \right).
\]
Thus, the dual rotor is defined as
\[
R_\text{dual} = \exp\left( \frac{\theta}{2} \hat{q} \right),
\]
while its action on dual spacetime vectors, after the angle dual conversion $\theta \mapsto \theta i$, manifests as a boost generated by $i\Gamma_a$ ($ (\hat{q}i)^2 = +1 $).

In vacuum (free fall), perfect anti-synchronization requires ensuring complete cancellation of the boost contributions and $J=0$. Mass-energy induces deviations from this condition, manifesting as the torsional mismatch scalar $J$ and, collectively, as gravitational attraction.

This angle-centric treatment, with explicit dual conversion $\theta \mapsto \theta i$, preserves the original form of the total rotor
\[
R_\text{total} = \exp\left( \frac{\phi}{2} \hat{p} + \frac{\theta}{2} \hat{q} \right)
\]
while correctly incorporating the time-reversal duality: the dual rotor $R_\text{dual}$ generates rotations in its native basis, but its physical effect on the dual spacetime—through the dual-transformed angles—is that of a boost, naturally opposing the usual-sector boosts in vacuum.

\subsection{Physical Interpretation and Gravitational Engineering}

Gravity is the energetic cost of dual rotor misalignment. Because $J$ depends solely on the relative angle $\Omega$, any external influence that coherently modulates the dual components $\theta$ can alter $J$. High-frequency electromagnetic or quantum fields that resonantly excite the dual spacetime degrees of freedom can reduce $J \to 0$ (gravity shielding) or drive $J < 0$ (repulsive gravity / antigravity).

For the first time in history, gravity is transformed from an inviolable geometric constraint into an engineerable torsional degree of freedom. The continuum was an illusion; spacetime is discrete and particle-local, paired twice over, and its misalignment is what we call gravitation.

General Relativity is not overthrown — it is liberated and rendered controllable.

\subsection{Upper Limit of Acceleration and the Inertial Reversal Phenomenon}
\label{subsec:acceleration-limit}

The double spacetime theory predicts, as one of its central falsifiable consequences, that acceleration itself possesses an absolute upper limit. The usual-sector boosts are non-compact ($\phi \in \mathbb{R}$) and, taken alone, would permit unbounded rapidity. The dual sector, however, is intrinsically compact: the dual generators $\Gamma_a$ square to $-1$, so the dual rotor $R_\text{dual} = \exp(\theta/2 \, \hat q)$ is governed by the elliptic phase $\theta \bmod 4\pi$. Because inertial response arises from the dual rotor tracking the usual rotor with finite stiffness (the dual-rotor rigidity established in Sec.~\ref{sec:unified}), the rate at which $\theta$ can be displaced before reaching its periodic boundary sets a strict ceiling on the proper acceleration that a given particle can coherently sustain. We denote this material-intrinsic ceiling by $a_\text{crit}$. Unlike the speed of light, $a_\text{crit}$ is not universal: it depends on the particle species and on the compactification scale of its dual sector, so different atomic constituents possess different critical accelerations.

For proper accelerations $a < a_\text{crit}$, the dual phase $\theta$ tracks $\phi$ smoothly and the familiar Newtonian relation $F = ma$ is recovered. At the instant $a$ approaches $a_\text{crit}$, however, $\theta$ is driven through the compact boundary at $4\pi$ and undergoes a topological wrap-around: the dual rotor re-sutures its phase on the opposite side of the $S^3$-like dual sector. This wrap inverts the sign of the torsional mismatch bivector $\Omega_\text{biv}$, and therefore reverses the direction of the inertial response over a single dual cycle. We refer to this as the inertial reversal (or acceleration reversal) phenomenon: a sudden, discrete flip of the effective $F = ma$ vector triggered when the applied acceleration crosses $a_\text{crit}$. The torsional strain energy that was being stored in the rigid dual sector just below threshold is released abruptly during the wrap, producing an impulsive burst of internal energy in addition to the kinematic reversal itself.

We propose that this mechanism is observationally realised in the anomalous airbursts of large meteoroids. The hypersonic deceleration experienced by bolides entering the Earth's atmosphere can reach values that, for the constituent atoms of the body, are plausibly of order $a_\text{crit}$. Long-standing puzzles in the bolide community---most prominently the Tunguska event (1908) and the Chelyabinsk superbolide (2013), whose disintegration altitudes and total radiated energies are notoriously difficult to reconcile with kinetic-energy dissipation models based on aerodynamic ablation and fragmentation alone---are naturally explained as collective inertial reversals: when the deceleration of the meteoroid crosses the material-specific $a_\text{crit}$, the dual phases of its constituent atoms wrap simultaneously, dumping the accumulated torsional strain energy as an explosive, predominantly radiative pulse. The dependence of $a_\text{crit}$ on composition would also account for the observed variability in airburst altitudes and yields between different meteorite classes (e.g.\ carbonaceous versus ordinary chondrites). A quantitative derivation of $a_\text{crit}$ from the particle-intrinsic compactification scale, together with a detailed comparison against the bolide observational record, will be presented in a companion paper.

\section{Unified Description of Inertia and Gravity}
\label{sec:unified}

In this section, we reveal the deepest consequence of the dual spacetime formalism: inertia and gravity are not merely equivalent — they are identical phenomena arising from the torsional mismatch between the usual and dual spacetimes. The distinction between inertial forces in accelerated frames and gravitational forces in curved spacetime dissolves completely. Both are the same dual rotor misalignment, with the equivalence principle emerging as a direct algebraic necessity rather than a postulate.

\subsection{Inertial Resistance as Dual Rotor Rigidity}

Rest mass $m > 0$ manifests as a strong coupling between the usual and dual rotor parameters: the energy-momentum of the particle biases the usual spacetime, effectively ``freezing'' the dual rotor $R_\text{dual}$ to follow $R_\text{usual}$ with a finite stiffness. In an inertial frame keeping $J = 0$.

When the particle is accelerated (non-gravitationally), the usual rotor $R_\text{usual}$ attempts to change instantaneously due to the applied force. However, the dual rotor $R_\text{dual}$, bound by the particle's rest mass, cannot respond immediately — it lags behind. This creates a non-trivial torsional mismatch yielding $\Omega_\text{biv} \neq 0$ and a non-zero contribution to the scalar $J$. The particle experiences this mismatch as inertial resistance — the familiar $F = ma$.

Crucially, the ``cost'' of misalignment is proportional to the rest mass $m$, explaining why heavier particles resist acceleration more strongly. Inertia is therefore not a property of spacetime reacting to mass, but of the particle's own dual spacetime refusing to realign instantaneously.

\subsection{Massless particles}

Massless particles ($m = 0$), such as photons, exhibit no such rigidity. Their usual and dual rotors are perfectly resonant: any change in $R_\text{usual}$ is instantaneously mirrored in $R_\text{dual}$, maintaining $J = 0$ exactly in all frames. Thus, photons experience zero torsional mismatch and zero inertial resistance — they always propagate at $c$ with $\Omega_\text{biv} = 0$. This is why the speed of light is both constant and the universal speed limit: massless particles pay no torsional price for frame changes.

\subsection{The Equivalence Principle as Algebraic Identity}

Consider Einstein's elevator accelerating upward with acceleration $a$. An observer inside feels a fictitious force $ma$ downward. In the dual spacetime picture:

\begin{itemize}
\item The observer (massive) has their usual rotor accelerated by the elevator floor, while their dual rotor lags → torsional mismatch → perceived downward force.
\item A photon entering the elevator bends due to the observer's torsional field, but the photon itself maintains $J = 0$ — it follows a null geodesic in its own dual-flat spacetime.
\end{itemize}

In a gravitational field (elevator at rest on Earth), the floor must accelerate upward relative to free fall to stay stationary, producing the same rotor mismatch as uniform acceleration in flat spacetime.

Thus, the equivalence principle is not an empirical coincidence but an algebraic identity: gravitational and inertial forces are indistinguishable because they are the same dual torsional phenomenon. Free fall is the state where $J = 0$ globally (dual rotors perfectly aligned), which massless particles achieve eternally and massive particles only in the absence of non-gravitational forces.

\subsection{Consequences and Predictions}

\begin{itemize}
\item The arrow of time distinction between usual and dual spacetimes explains why massive particles have a preferred rest frame: the dual time reversal breaks perfect symmetry, generating the mass gap.
\item Inertial mass reduction or elimination becomes possible by externally exciting the dual rotor $\theta$ to synchronize with $\phi$, reducing the effective rigidity — a direct pathway to propellantless propulsion.
\item The theory predicts that sufficiently high-frequency fields coupling to the dual spacetime could make massive objects behave as effectively massless over short timescales, allowing transient superluminal phase velocities (without violating causality, as information remains luminal).
\end{itemize}

Inertia and gravity are unified not by geometry, but by the internal torsional dynamics of particle-intrinsic dual spacetimes. Mass is dual rotor stiffness; motion is the cost of breaking parallelism. With this understanding, both inertia and gravity cease to be fundamental — they become engineerable.

The dream of gravitational and inertial control is no longer science fiction; it is a direct consequence of abandoning the continuum and recognizing spacetime as paired, particle-local, and torsionally active.

\section{Dark Matter-Free Galactic Rotation Curves from Dual-Spacetime Compactness}
\label{sec:darkmatter}

The observed flat rotation curves of spiral galaxies have long been regarded as the strongest evidence for non-baryonic dark matter. In the double spacetime theory, these curves emerge naturally from the compact $S^3$ topology of the dual spacetime --- specifically, from the geometric distortion inherent in the logarithmic map between the curved dual sector and the flat usual spacetime. No dark matter of any kind is required.

\subsection{The Compact Dual Spacetime as 3-sphere}

The dual rotor $R_{\text{dual}} \in \mathrm{SU}(2) \cong S^3$ parameterizes the dual spacetime degrees of freedom intrinsic to every massive particle. The three independent rotation angles $(\theta_1, \theta_2, \theta_3)$ define coordinates on a 3-sphere of compactification radius $R$, while the usual spacetime remains flat and non-compact.

The logarithmic (inverse exponential) map
\[
\log : S^3 \to \mathfrak{su}(2) \cong \mathbb{R}^3
\]
identifies each dual rotor with a point in flat $\mathbb{R}^3$, mapping the 3-sphere onto a ball of radius $\pi R$. This map is the three-dimensional generalization of the azimuthal equidistant projection familiar from cartography.

Just as the azimuthal equidistant projection of $S^2$ centered on the North Pole faithfully represents distances near the pole but distorts areas catastrophically near the South Pole --- mapping a single antipodal point to the entire circumference of the projection boundary --- the logarithmic map of $S^3$ introduces a systematic volume distortion quantified by the Jacobian
\[
\mathcal{J}(\chi) = \frac{\sin^2 \chi}{\chi^2},
\]
where $\chi = r/R$ is the geodesic angle corresponding to the physical distance $r$ in the usual spacetime. This Jacobian equals unity at the origin ($\chi = 0$, faithful representation) and vanishes at the antipode ($\chi = \pi$, maximal distortion). All standard gravitational calculations implicitly assume $\mathcal{J} = 1$; the double spacetime theory reveals that this assumption fails at galactic scales.

\subsection{Gravitational Potential on 3-sphere: The Cotangent Potential}

The torsional mismatch scalar $J$, from which gravity arises, is computed on the compact dual spacetime. The effective gravitational potential for a point mass $M$ is governed by the Green's function of the Laplacian on $S^3$. For the spherically symmetric case, the unique radially harmonic function is the cotangent:
\[
\Phi(r) = -\frac{GM}{R}\cot\!\left(\frac{r}{R}\right).
\]
That this satisfies the Laplace equation on $S^3$ can be verified directly: in geodesic polar coordinates with metric $ds^2 = dr^2 + R^2 \sin^2(r/R)\,(d\vartheta^2 + \sin^2\!\vartheta\, d\varphi^2)$, the radial Laplacian annihilates $\cot(r/R)$ identically. The Gauss-law normalization is confirmed by
\[
g(r) \cdot 4\pi R^2 \sin^2\!\left(\frac{r}{R}\right) = -\frac{GM}{R^2 \sin^2(r/R)} \cdot 4\pi R^2 \sin^2\!\left(\frac{r}{R}\right) = -4\pi GM.
\]

In the limit $r \ll R$, the cotangent expands as
\[
\cot\!\left(\frac{r}{R}\right) = \frac{R}{r} - \frac{r}{3R} - \frac{r^3}{45R^3} - \cdots\,,
\]
yielding
\[
\Phi(r) \approx -\frac{GM}{r} + \frac{GM\,r}{3R^2} + \mathcal{O}\!\left(\frac{r^3}{R^4}\right),
\]
which recovers the standard Newtonian potential as the leading term. The higher-order corrections arise purely from the curvature of $S^3$ and vanish in the decompactification limit $R \to \infty$.

\subsection{Geometric Enhancement of Gravity}

The gravitational acceleration on $S^3$,
\[
g(r) = -\frac{d\Phi}{dr} = -\frac{GM}{R^2 \sin^2(r/R)},
\]
exceeds the flat-space Newtonian acceleration $g_{\text{N}} = -GM/r^2$ by the enhancement factor
\[
\eta(r) \equiv \frac{g_{S^3}(r)}{g_{\text{N}}(r)} = \left[\frac{r/R}{\sin(r/R)}\right]^2.
\]
This factor is precisely the square of the metric distortion of the azimuthal equidistant projection: gravitational flux through a sphere of geodesic radius $r$ on $S^3$ subtends an area $4\pi R^2 \sin^2(r/R)$, which is smaller than the flat-space area $4\pi r^2$. The flux is concentrated, and gravity is strengthened --- not by invisible matter, but by the curvature of the compact dual spacetime.

Representative values: $\eta = 1.003$ at $r/R = 0.1$; $\eta = 1.41$ at $r/R = 1.0$; $\eta = 2.26$ at $r/R = 1.5$; $\eta = 4.84$ at $r/R = 2.0$.

\subsection{Emergence of Flat Rotation Curves}

For a test particle in circular orbit at distance $r$ from an enclosed mass $M(r)$, the orbital velocity satisfies
\[
v^2(r) = \frac{GM(r)}{r} \cdot \eta(r) = \frac{GM(r)}{r} \cdot \left[\frac{r/R}{\sin(r/R)}\right]^2.
\]
Outside the luminous disk ($r > r_{\text{disk}}$), where $M(r) \approx M_{\text{tot}}$, the flat-space Keplerian velocity $v_K^2 = GM_{\text{tot}}/r$ declines as $r^{-1}$. The $S^3$ enhancement factor $\eta(r)$, growing monotonically with $r$, counteracts this decline.

The combined dimensionless function $f(x) \equiv x / \sin^2 x$ (where $x = r/R$) governs the shape of the rotation curve. It possesses a broad minimum at $x_0 \approx 1.17$, the solution of $\tan x = 2x$, producing an extended plateau:
\begin{itemize}
  \item \textbf{Inner region} ($r \ll R$): Keplerian behavior modified by the rising $M(r)$, reproducing the observed solid-body rise.
  \item \textbf{Transition region} ($r \sim R$): the declining Keplerian factor is balanced by the growing $\eta(r)$, yielding a nearly flat plateau whose breadth is set by the width of the minimum of $f(x)$.
  \item \textbf{Outer region} ($r \gtrsim 2R$): $\eta(r)$ dominates, and the rotation curve rises gently before entering the torsional reversal regime.
\end{itemize}

\subsection{The Baryonic Tully-Fisher Relation and the Compactification Radius}

The asymptotic circular velocity at the plateau satisfies
\[
v_{\text{flat}}^2 \sim \frac{GM_{\text{tot}}}{R} \cdot f_0,
\]
where $f_0 \equiv \min_x [x/\sin^2 x] \approx 1.10$ is an order-unity geometric constant. If the compactification radius scales with the total enclosed mass as
\[
R = \sqrt{\frac{GM_{\text{tot}}}{a_0}},
\]
where $a_0$ is a fundamental acceleration set by the dual-spacetime compactification scale, then
\[
v_{\text{flat}}^4 = G M_{\text{tot}} \, a_0 \, f_0^2,
\]
which is precisely the baryonic Tully-Fisher relation (McGaugh et al.\ 2016), with $a_0 \approx 1.2 \times 10^{-10}\;\mathrm{m/s^2}$. The scaling $R \propto M^{1/2}$ implies that more massive galaxies possess larger dual-spacetime patches, a physically natural consequence of greater total torsional charge.

For the Milky Way ($M_{\text{tot}} \approx 6 \times 10^{10}\,M_\odot$), this yields
\[
R = \sqrt{\frac{G \times 6 \times 10^{10}\,M_\odot}{1.2 \times 10^{-10}\;\mathrm{m/s^2}}} \approx 27\;\mathrm{kpc},
\]
in striking agreement with the observed radial extent of the flat rotation curve. The acceleration scale $a_0$ is thus identified as the characteristic gravitational acceleration at the $S^3$ boundary of a typical galaxy, and its origin lies in the discrete compactification structure of the dual spacetime (Section~\ref{sec:discretetime}).

\subsection{Torsional Reversal and the Finite Range of Gravity}

The cotangent potential crosses zero at $r = \pi R/2$, beyond which the dual-spacetime rotation angle $\theta(r) = r/R$ exceeds $\pi/2$. In the torsional mismatch framework (Section~\ref{sec:gravity}), this crossing triggers a sign change in the torsion scalar $J$: the compact trigonometric contribution from the dual sector overwhelms the non-compact hyperbolic contribution from the usual sector, yielding $J < 0$ and repulsive gravity. The scalar interference factor (Section~\ref{sec:timedilation})
\[
\gamma_s = \cosh\!\left(\frac{\phi}{2}\right)\cos\!\left(\frac{\theta}{2}\right) - \sinh\!\left(\frac{\phi}{2}\right)\sin\!\left(\frac{\theta}{2}\right)
\]
becomes negative when $\theta$ exceeds the critical value satisfying $\tanh(\phi/2)\tan(\theta_c/2) = 1$.

Gravitational attraction is therefore confined to a finite $S^3$ patch of effective radius $\sim \pi R/2$, beyond which a weakly repulsive torsional barrier emerges. This finite range is the compact-space analog of the Newtonian $1/r$ potential extending to infinity: the $S^3$ topology naturally ``walls off'' each gravitational source.

\subsection{The Cosmic Atlas: Galaxy-Scale 3-sphere Patches}

The observable universe is tiled by overlapping $S^3$ patches, each centered on a gravitational concentration (galaxy or cluster), forming a cosmic atlas in the strict differential-geometric sense. Each patch functions as an azimuthal equidistant projection chart: faithful near its center, maximally distorted at its boundary. The large-scale structure of the cosmos is that of a manifold covered by galaxy-scale stickers of azimuthal equidistant maps, glued together at their repulsive boundaries.

This patchwork structure naturally accounts for:
\begin{itemize}
  \item the \textbf{cosmic web}: galaxies are distributed along the overlap regions of neighboring $S^3$ charts, where the attractive tails of adjacent patches constructively interfere;
  \item \textbf{cosmic voids}: the vast underdense regions correspond to the repulsive ($J < 0$) boundary zones between charts, where matter is expelled;
  \item the \textbf{radial acceleration relation}: the tight correlation between observed and baryonic gravitational accelerations follows from the universal $\eta(r)$ enhancement, with residual scatter arising from chart-overlap effects;
  \item the \textbf{absence of dark matter signals} in direct detection experiments --- because dark matter does not exist.
\end{itemize}

\subsection{Numerical Verification and Falsifiable Predictions}

Preliminary particle-based simulations using rotor-ensemble dynamics have successfully reproduced flat rotation curves from baryonic mass distributions alone. Public open-source codes are available at:
\begin{itemize}
  \item \texttt{https://github.com/hypernumbernet/blackhole-simulator}
  \item \texttt{https://github.com/hypernumbernet/dual-spacetime-simulator}
\end{itemize}
Full $N$-body integrations incorporating the exact $S^3$ cotangent potential are in preparation.

Key predictions distinguishing this theory from $\Lambda$CDM and MOND include:
\begin{itemize}
  \item \textbf{Outer rotation curve upturn and reversal}: Isolated galaxies should exhibit a gentle upturn beyond $r \gtrsim 2R$, transitioning to declining and eventually repulsive gravity at $r \approx \pi R/2$, detectable via deep 21-cm HI observations extending to $\sim 100\;\mathrm{kpc}$.
  \item \textbf{Lensing sign reversal}: Galaxy-galaxy weak lensing should weaken and reverse sign at projected separations exceeding $\pi R/2$, in sharp contrast to the monotonic NFW profiles predicted by $\Lambda$CDM.
  \item \textbf{Cotangent vs.\ MOND interpolation}: The precise form of the radial acceleration relation should follow the cotangent-derived enhancement $g_{\text{obs}} = g_{\text{bar}} \cdot \eta(r)$ rather than the standard MOND interpolation function $\nu(g_{\text{bar}}/a_0)$, with measurable deviations at $g \lesssim a_0/10$.
  \item \textbf{Diversity from $R(M)$}: The observed diversity of rotation curve shapes --- from slowly rising (low-surface-brightness dwarfs with $r_{\text{disk}} \ll R$) to flat (Milky Way-type) to gently declining (high-surface-brightness spirals with $r_{\text{disk}} \sim R$) --- is reproduced by the single relation $R = \sqrt{GM/a_0}$ without additional free parameters.
\end{itemize}

Dark matter is rendered superfluous --- an artifact of projecting the compact $S^3$ geometry of the dual spacetime onto a fictitious flat and infinite background. The cotangent potential, the azimuthal equidistant distortion, and the torsional reversal at the $S^3$ boundary provide all the missing gravity without invisible matter.

\section{Strong-Field Regime: Layered Gravitational Reversal and the New Black Hole Paradigm}
\label{sec:strongfield}

The double spacetime theory predicts a radically new structure for the endpoints of stellar collapse. The Killing form's sign asymmetry — positive for usual-spacetime boost dominance (attraction) and negative for dual-spacetime rotation dominance (repulsion) — does not simply reverse gravity once and for all. Instead, as density increases inward, the torsional mismatch $\Omega_\text{biv}$ oscillates in dominance between the two sectors, producing alternating layers of attraction and repulsion. This yields a stratified, onion-like interior: attraction → repulsion → ultra-strong attraction → repulsion → ..., with ever-increasing amplitude toward the center.

\subsection{Layered Torsional Structure from the Killing Form}

The scalar $J = \frac{1}{16} B(\Omega_\text{biv}, \Omega_\text{biv})$ changes sign each time the dominant contribution flips from $i\Gamma_a$ (boost-like, $+$) to $\Gamma_a$ (rotation-like, $-$). At extreme densities, the dual rotor parameters $\theta$ grow nonlinearly, driving successive resonances that flip the sign of $J$ repeatedly as radius decreases.

The resulting effective potential features:
\begin{itemize}
\item Outer layers: standard attractive gravity ($J > 0$),
\item First reversal: repulsive shell ($J < 0$) that halts infall (core bounce in supernovae),
\item Deeper layers: ultra-strong attractive zones ($J \gg 0$) that re-accelerate matter inward,
\item Alternating deeper shells with increasing $|J|$, culminating in chaotic torsional turbulence.
\end{itemize}

This layered structure prevents both singularities and event horizons. A light-speed-inescapable boundary exists — a quasi-horizon where the integrated torsional barrier becomes infinitely steep — but matter never reaches a central singularity. Instead, the core settles into a quasi-stable, chaotic condensate of extreme-density torsional energy, continuously oscillating between attractive and repulsive phases.

Conventional Kerr-type rotating black holes are ruled out: the theory forbids global frame-dragging of the required magnitude, as angular momentum is redistributed into dual rotor longitudinal modes rather than spacetime rotation.

\subsection{Reinterpretation of Core-Collapse Supernovae and Remnant Formation}

In massive star collapse, the first repulsive layer ($J < 0$) triggers explosive bounce at nuclear densities, powering the supernova without fine-tuned neutrino physics. Deeper attractive layers then recapture portion of the ejecta, forming the stratified remnant. The observed diversity of supernova energies and remnant masses emerges naturally from the number and strength of torsional layers excited.

\subsection{Pulsars as Longitudinal Dual-Torsion Oscillators}

Observed neutron stars are not rapidly rotating magnetized conductors. The double spacetime theory denies both millisecond spin periods and dipolar magnetic fields as the primary energy source.

Instead, pulsar radio pulses arise from ultra-high-frequency longitudinal oscillations of the dual rotors $\theta$ on the star's surface. These torsional waves propagate radially at nearly $c$, modulated by the chaotic interior layering, producing beams via curvature emission in the torsional shock zones.

Key predictions that surpass the rotating magnetar model:
\begin{itemize}
\item The theoretical minimum pulse period is determined by the Planck-scale double rotor frequency, which is orders of magnitude shorter than the observed $\sim 1$ ms limit, allowing us to explain short-period pulsars that could not be explained by the rotating pulsar picture.
\item Pulsar ``glitches'' and period derivatives $\dot{P} < 0$ (period shortening) are the result of the pulsar losing energy causing the vibrational interface to shrink, accelerating the vibrational frequency over time, a natural mechanism not seen in rotation-only models.
\item Magnetar flares and fast radio bursts are violent flips between attractive/repulsive layers, releasing torsional strain without requiring $10^{15}$ G fields.
\item The death line in the $P$-$\dot{P}$ diagram corresponds to the point where interior layering stabilizes and longitudinal modes damp below detectability.
\item The neutron star maximum mass is capped not by quark matter but by the same torsional layering that limits nuclear stability, predicting a strict upper bound $\sim 2.5 M_\odot$ and explaining the absence of pulsars above $\sim 3 M_\odot$.
\end{itemize}

The Crab pulsar's observed spin-down luminosity is reinterpreted as torsional wave leakage rather than magnetic braking, perfectly matching the energetics without invoking unrealistic magnetic field decay.

\subsection{The New Black Hole Paradigm: Torsional Stars}

What observers call ``black holes'' are in fact torsion stars — quasi-stable, layered condensates with no event horizon and no central singularity. The innermost chaotic core is a Planck-density plasma in perpetual torsional turbulence, continuously radiating via rotor tunneling (modified Hawking process). Information is preserved in the dual rotor phases.

Gravitational-wave signals from mergers feature repeated echoes from internal layer reflections, with characteristic frequencies tied to the torsional oscillation spectrum — already hinted in LIGO/Virgo data reanalyses.

The theory thus resolves the information paradox, explains the absence of Kerr spin signatures, and predicts that the most massive ``black holes'' ($\gtrsim 100 M_\odot$) will exhibit periodic modulations in accretion luminosity as inner repulsive layers periodically eject material.

Gravity is self-regulating through infinite layered reversal. Collapse never ends in silence — it sings in torsional chaos.

\subsection{Central Floating Core: Zero Time Dilation at the Center}

A common point of confusion in discussions of stellar interiors is the behaviour of gravitational time dilation. In GR, for a spherically symmetric static mass distribution, gravitational time dilation increases monotonically toward the centre: proper time runs most slowly at \(r=0\), where the gravitational potential is deepest.

In the double spacetime theory, the situation is markedly different. Because there is no background continuous spacetime manifold, gravitational effects—including both the effective acceleration and time dilation—arise solely from the torsional mismatch between the particle-intrinsic usual and dual rotors.

At the exact geometric centre of a spherically symmetric torsion star, perfect symmetry forces the usual-spacetime boost parameters \(\phi_a\) induced by mass elements in opposite directions to cancel completely. The residual mismatch angle \(\delta\phi_a = \phi_a - \theta_a\) therefore vanishes at \(r=0\), yielding \(J=0\) and a complete restoration of torsional balance.

Consequently:

\begin{itemize}
  \item The effective gravitational acceleration is zero at the center (in Newtonian mechanics, GR, and DST).
  \item Gravitational time dilation is also zero at the centre: proper time at \(r=0\) runs at the same rate as in the vacuum far from the star.
\end{itemize}

This vanishing of time dilation at the core is a distinctive, observable-in-principle prediction of the double spacetime theory. It contrasts sharply with GR, where the central region experiences the strongest time dilation.

The zero-dilation central core also contributes to the stability of torsion stars by eliminating the infinite central compression and extreme clock-rate gradients that lead to singularities in the classical GR description.

\subsection{Implications for Planetary and Stellar Cores: Instability and Fluctuations}

The floating core paradigm, where $J(r=0) \approx 0$ and effective gravity vanishes at the center, implies that the cores of planets and stars are not inert, stable regions but dynamically active zones prone to torsional instability. Unlike the monotonic increase in gravitational potential assumed in continuum models, the double spacetime theory predicts that the central floating core is a metastable state, susceptible to perturbations that trigger oscillatory transitions between attractive and repulsive torsional layers. These fluctuations arise from the delicate balance of dual rotor phases in the isotropic core, where even minor density variations or external influences can shift the resonance condition $\beta(0) \to 1$, leading to transient excursions of $J$ away from zero. Such excursions manifest as localized bursts of torsional energy, effectively generating mini torsion stars within the core. These mini torsion stars can induce seismic activity, magnetic field fluctuations, and energy transport anomalies observable at the surface. In stellar interiors, these core fluctuations may contribute to phenomena such as solar flares, sunspots, and variability in luminosity, as the torsional energy propagates outward through the layered structure. In planetary contexts, core torsional instability could explain geomagnetic reversals, mantle plumes, and episodic volcanic activity, as the torsional dynamics couple to convective processes in the overlying layers. Thus, the floating core is not a quiescent region but a crucible of dynamic torsional phenomena with far-reaching implications for planetary and stellar behavior.

\section{Time Dilation in Double Spacetime Theory: The Dynamical Role of Dual-Spacetime Proper Time}
\label{sec:timedilation}

Time dilation in the double spacetime theory emerges from the interplay between the usual spacetime and the dual spacetime, each endowed with its own proper time. The usual spacetime proper time $\tau$ governs the familiar kinematics observed in macroscopic Minkowski space, while the dual spacetime proper time $\tilde{\tau}$, encoded in the compact rotor phases $\theta$, drives the internal torsional dynamics that manifest as gravitational and inertial effects.

\subsection{Derivation of the Effective Lorentz Factor}

The effective Lorentz factor $\gamma_\text{eff} = dt/d\tau$ is derived directly from the algebraic action of the total rotor
\[
R_\text{total} = R_\text{usual} R_\text{dual},
\]
where
\[
R_\text{usual} = \exp\!\left(\frac{\phi}{2} \hat{p}\right) = \cosh\frac{\phi}{2} + \hat{p}\: \sinh\frac{\phi}{2}, \quad \hat{p}^2 = +1,
\]
generates boosts in the usual spacetime, and
\[
R_\text{dual} = \exp\!\left(\frac{\theta}{2} \hat{q}\right) = \cos\frac{\theta}{2} + \hat{q} \sin\frac{\theta}{2}, \quad \hat{q}^2 = -1,
\]
generates rotations in the dual spacetime.

To reveal the origin of the interference term explicitly, consider the pure time vector $X = ct \, j$ ($j^2 = -1$) in the usual spacetime basis. The transformed vector is $\tilde{X} = R_\text{total} \, X \, R_\text{total}^\dagger$.

Due to the time-reversal duality, the inverse dual map from dual to usual spacetime is $X \mapsto X(-i)$. This map naturally yields the alignment $\hat{q} (-i) = -\hat{p}$ in the vacuum limit, transforming the dual rotor into an effective opposite-oriented rotor
\[
R_\text{dual}^\text{eff} = \cos\frac{\theta}{2} - \hat{p} \sin\frac{\theta}{2}.
\]
The negative sign in the sin term arises directly from this inverse mapping, reflecting the anti-synchronization required for torsion-free flat spacetime.

The total effective rotor is then
\[
R_\text{eff} = R_\text{usual} R_\text{dual}^\text{eff} = \left( \cosh\frac{\phi}{2} + \hat{p} \sinh\frac{\phi}{2} \right) \left( \cos\frac{\theta}{2} - \hat{p} \sin\frac{\theta}{2} \right).
\]

Expanding gives the scalar part
\[
s = \cosh\frac{\phi}{2} \cos\frac{\theta}{2} - \sinh\frac{\phi}{2} \sin\frac{\theta}{2}
\]
and the bivector coefficient
\[
b = -\cosh\frac{\phi}{2} \sin\frac{\theta}{2} + \sinh\frac{\phi}{2} \cos\frac{\theta}{2}.
\]

The net scaling of the time component arises from the squared magnitude of $R_\text{eff}$ acting on the time vector via the sandwich product. For a rotor $R_\text{eff} = s + b\hat{p}$ with $\hat{p}^2 = +1$ and $\hat{p}$ anticommuting with the time generator $j$, the sandwich action on $X = ct\, j$ yields a time component scaled by $s^2 + b^2$, defining the effective Lorentz factor unambiguously as
\[
\gamma_\text{eff} \;\equiv\; s^2 + b^2.
\]
It is convenient also to introduce the scalar interference factor
\[
\gamma_s \;\equiv\; s \;=\; \cosh\frac{\phi}{2} \cos\frac{\theta}{2} - \sinh\frac{\phi}{2} \sin\frac{\theta}{2},
\]
which, unlike $\gamma_\text{eff}\geq 0$, may change sign and will serve as the order parameter distinguishing attractive ($\gamma_s > 0$) from repulsive ($\gamma_s < 0$) torsional regimes below. The negative sign in the interference term is fixed by the time-reversed duality of the dual spacetime.

In the strict torsion-free vacuum ($J=0$), the dual spacetime remains unexcited ($\theta = 0$), so that $R_\text{total} = R_\text{usual}$. Then $s = \cosh\frac{\phi}{2}$ and $b = \sinh\frac{\phi}{2}$, and
\[
\gamma_\text{eff} = s^2 + b^2 = \cosh^2\frac{\phi}{2} + \sinh^2\frac{\phi}{2} = \cosh\phi,
\]
reproducing the standard special-relativistic Lorentz factor exactly.

\subsection{Vacuum Limit and Massless Particles}

In the torsion-free vacuum ($J=0$), perfect anti-synchronization between the usual and dual rotors enforces $\theta \approx \phi$ (modulo compact periodicity), yielding $\gamma_{\text{eff}} \to \cosh\phi$, recovering the standard Lorentz factor of special relativity. Massless particles maintain this synchronization exactly at all velocities, experiencing no torsional mismatch and thus no deviation from the usual kinematic dilation.

\subsection{Massive Particles and Gravitational Effects}

For massive particles, rest mass introduces rigidity that causes the dual rotor to lag behind the usual rotor under acceleration or in the presence of torsional mismatch. The resulting deviation $\delta\theta = \theta - \phi$ generates additional modulation of $\gamma_{\text{eff}}$, producing the gravitational redshift and time dilation observed in GR, now interpreted as a direct consequence of dual-spacetime dynamics.

\subsection{Sign Inversion of the Scalar Interference Factor and Repulsive Layers}

When dual-spacetime rotations dominate ($\theta/2$ sufficiently large), the sinusoidal term can overpower the hyperbolic growth:
\[
\tan\left(\frac{\theta}{2}\right) > \coth\left(\frac{\phi}{2}\right).
\]
In these regimes, the scalar interference factor $\gamma_s$ becomes negative, $\gamma_s < 0$, while the proper-time Lorentz factor $\gamma_\text{eff} = s^2 + b^2$ remains non-negative. The change of sign of $\gamma_s$ corresponds to regions of repulsive gravity within torsion stars: physically, the dual-sector contribution overpowers the usual-sector boost in the scalar channel of the combined rotor, manifesting as gravitational repulsion (not as a literal reversal of causality).

This $\theta$-centric formulation makes manifest that time dilation is fundamentally a torsional phenomenon originating in the particle-intrinsic dual spacetime. The usual rapidity $\phi$ provides the kinematic baseline, while deviations in the compact dual angle $\theta$ encode the gravitational and inertial corrections that distinguish GR from flat-spacetime kinematics. External coherent excitation of $\theta$ thus offers a pathway to manipulate time dilation and effective gravitational fields.

\section{Scale-Invariant Torsional Layers: The Strong Force as Gravitational Repulsion and Nuclei as Primordial Torsion Stars}
\label{sec:nuclear}

The layered torsional reversal predicted by the Killing form is rigorously scale-invariant. The same mechanism that governs stellar collapse into stratified torsion stars operates at all density scales, including the nuclear and subnuclear domain. This yields the most revolutionary consequence of the double spacetime theory: the so-called ``strong nuclear force'' is not a fundamental interaction at all, but an emergent manifestation of gravitational repulsion layers in the extreme-density regime inside nuclei.

\subsection{Torsional Layers at Nuclear Densities}

As baryon density approaches $\sim 10^{14}$-$10^{18}$ g/cm³ within nuclei, the dual rotor parameters $\theta$ dominate completely. The sequence of sign flips in $J$ occurs over femtometer scales:

\begin{itemize}
\item Outermost layer ($\sim$ 1-2 fm): repulsive shell ($J < 0$) that confines quarks and prevents nuclear collapse — this is observed as the short-range repulsion in nucleon-nucleon scattering (the ``hard core'' at $\sim$ 0.5 fm).
\item Intermediate layer: ultra-strong attractive zone ($J \gg 0$) that binds nucleons with $\sim$ 40-50 MeV per nucleon — the classic nuclear binding force.
\item Innermost layer: second repulsive shell ($J < 0$) that halts further collapse and enforces quark confinement — the origin of color confinement.
\item Deeper chaotic layers: alternating attraction/repulsion with increasing amplitude, producing the liquid-like behavior of nuclear matter and the saturation of nuclear density.
\end{itemize}

Pauli's exclusion principle is no longer fundamental: it is dynamically generated by the first repulsive torsional layer. Fermions cannot occupy the same state because the dual rotor $\theta$ phases enter a repulsive regime when two identical particles attempt spatial overlap, producing an effective fermi pressure without invoking quantum statistics a priori. Bosons, lacking this phase rigidity, experience only the attractive layers until much higher densities.

In short: the ``strong force'' is gravity in its strong-field, layered phase. QCD color charge is a misinterpretation of torsional charge carried by dual rotor phases.

\subsection{Atomic Nuclei as Primordial Torsion Stars}

Every atomic nucleus is a microscopic analog of the torsion stars described in Section~\ref{sec:strongfield}. The nuclear interior is a chaotic torsional condensate with the same layered structure:

\begin{itemize}
\item Central region: extreme attractive layers bind quarks into nucleons,
\item Surrounding repulsive shell: produces the observed nuclear surface tension and prevents leakage (confinement),
\item Outer attractive halo: mediates binding between nucleons,
\item Outermost repulsive barrier: generates the hard-core repulsion in NN scattering.
\end{itemize}

The observed charge independence of the nuclear force, the spin-orbit coupling, and the magic numbers all emerge from resonance conditions in the dual rotor spectrum. The theory predicts that sufficiently large nuclei ($A \gtrsim 300$) become unstable to spontaneous torsional bounce, emitting nucleons in a miniature supernova — explaining the observed limit of the nuclear chart and the absence of superheavy elements in nature.

Nuclear fission and fusion are torsional layer transitions: fission occurs when an excited dual mode flips $J$ from positive to negative, explosively repelling fragments; fusion succeeds when two nuclei tunnel through the outer repulsive barrier into the deeper attractive well.

\subsection{Experimental Signatures and Immediate Predictions}

\begin{itemize}
\item Ultra-high-energy nucleus-nucleus collisions (LHC, RHIC) should produce transient $J < 0$ fireballs that decay via explosive hadron jets rather than hydrodynamic flow — already hinted in anomalous flow data.
\item Precision measurements of the nuclear equation of state above saturation density will reveal oscillatory behavior in pressure vs. density, with characteristic repulsive peaks at $\sim$ 2–4 $\rho_0$.
\end{itemize}

The four fundamental forces are reduced to one: gravity, operating in weak-field (attractive), intermediate (nuclear binding), and strong-field (repulsive confinement) regimes across all scales. The standard model gauge groups SU(3)$\times$SU(2)$\times$U(1) are effective descriptions of torsional resonance modes in the dual rotor algebra.

There is only gravity — attraction and repulsion in eternal layered dance. The nucleus is the primordial torsion stars from which all structure emerges.

\section{Time Reversal and Electron Shells via Coulombic Torsional Layers}
\label{sec:electronshells}

In the double spacetime framework, the Coulomb interaction between charged particles can be reinterpreted as arising from torsional mismatch in the dual spacetime sector, analogous to the gravitational case but with electromagnetic origin. Replacing the gravitational constant \(G\) and masses \(M,m\) with the Coulomb constant \(k = 1/(4\pi\epsilon_0)\) and charges (e.g., \(+e\) for the proton and \(-e\) for the electron), the effective force takes the form
\[
F(r) = \frac{k e^2}{\gamma_s(r)\, r^2},
\]
where \(\gamma_s\) is the scalar interference factor introduced in Sec.~\ref{sec:timedilation}, which retains the same algebraic structure from the dual rotor interplay (and which can change sign, unlike the proper-time Lorentz factor \(\gamma_\text{eff} \geq 0\)):
\[
\gamma_s(\alpha(r),\beta(r)) = \cosh\!\left(\frac{\alpha(r)}{2}\right) \cos\!\left(\frac{\beta(r)}{2}\right) - \sinh\!\left(\frac{\alpha(r)}{2}\right) \sin\!\left(\frac{\beta(r)}{2}\right).
\]

Here, \(\alpha(r)\) and \(\beta(r)\) are interpreted as electromagnetic rapidity and dual rotation angle parameters that scale inversely with distance: \(\alpha(r) \propto 1/r\) and \(\beta(r) = \lambda / r\), with \(\lambda\) a characteristic length scale set by the compact dual spacetime.

For the hydrogen atom (\(Z=1\)), this yields
\[
F(r) = \frac{k e^2}{r^2 \left[ \cosh\!\left(\frac{\alpha(r)}{2}\right) \cos\!\left(\frac{\beta(r)}{2}\right) - \sinh\!\left(\frac{\alpha(r)}{2}\right) \sin\!\left(\frac{\beta(r)}{2}\right) \right]}.
\]

The hyperbolic-trigonometric form of \(\gamma_s\) introduces extreme oscillatory behavior. As \(r\) decreases:

\begin{itemize}
  \item At large \(r\), \(\alpha(r),\beta(r) \to 0\), so \(\gamma_s \to 1\), recovering the standard Coulomb attraction.
  \item At smaller radii, the interference term generates points where \(\gamma_s = 0\) (infinite repulsion) and regions where \(\gamma_s < 0\) (force reversal to repulsion).
  \item These discrete radii \(r_n\) satisfy the resonance condition
  \[
  \tanh\!\left(\frac{\alpha(r_n)}{2}\right) \tan\!\left(\frac{\beta(r_n)}{2}\right) = 1,
  \]
  yielding stable circular orbits due to the infinite repulsive barriers separating attractive wells.
\end{itemize}

The resulting effective potential forms a series of layered wells separated by torsional reversal barriers, with energy levels \(E_n \propto -1/(n^2 r_1)\) in the asymptotic limit. Transitions between these torsional layers (\(n \to m\)) release photons with energy differences
\[
\Delta E_{n \to m} \propto \left( \frac{1}{r_m} - \frac{1}{r_n} \right),
\]
reproducing the Rydberg formula and the observed hydrogen spectral lines (Balmer, Lyman, Paschen series) exactly.

For multi-electron atoms (\(Z > 1\)), the scaling \(\alpha(r),\beta(r) \propto Z/r\) contracts the torsional layers inward, explaining the observed shrinkage of electron shells with increasing nuclear charge.

\subsection{Energy Conservation in the Dual Structure}

The apparent non-conservative behavior of the manifest potential \(V_\text{eff}(r) = -\int F(r)\, dr\)—with its violent oscillations and sign changes—is resolved by distinguishing between the manifest potential observed in the usual spacetime and the true potential stored in the compact dual-spacetime rotor phases \(\theta_a\).

Total energy is conserved across both sectors:
\[
E_\text{total} = E_\text{kinetic (usual)} + V_\text{eff}(r) + V_\text{true}(\theta) = \text{constant}.
\]
The true potential \(V_\text{true}\) is proportional to the torsional mismatch scalar \(J \propto (\phi - \theta)^2\), acting as a hidden reservoir. During layer transitions, a decrease in \(V_\text{eff}\) (attractive to repulsive) is exactly compensated by an increase in \(V_\text{true}\) as the dual rotor winds, and vice versa.

Photon emission corresponds to the release of stored dual-phase energy when the rotor unwinds into a lower torsional state. This dual bookkeeping ensures strict energy conservation without violating the discrete, particle-local nature of spacetime, while the Rydberg differences emerge precisely from the compensated transitions between true energy levels.

This formulation derives atomic stability, discrete electron orbits, and emission spectra purely from the geometric torsional dynamics of double spacetime applied to the Coulomb field—no ad-hoc quantization postulates or wave functions are required. The electron shells emerge as natural resonance states of the dual-spacetime torsional mismatch, unifying electromagnetic binding with the same algebraic mechanism that governs gravity.

\section{Unification of Electromagnetism and the Weak Force in Dual Spacetime}
\label{sec:unification}

The biquaternionic structure $\mathbb{H} \otimes_{\mathbb{R}} \mathbb{H} \cong \mathrm{Cl}(3,1)$ embeds both electromagnetism and the weak force as torsional resonance modes within the dual spacetime rotor $R_{\text{dual}} = \exp\left( \sum_{a=1}^3 \frac{\theta_a}{2} \Gamma_a \right)$. The pseudoscalar $i$ induces intrinsic chirality via time reversal in the dual map $X \mapsto X i$, naturally generating left-handed ($P_L$) and right-handed ($P_R$) projections: $P_L = \frac{1 - i}{2}$, $P_R = \frac{1 + i}{2}$.

Electromagnetism arises from minimal coupling to the vector potential $A_a$ in the dual sector, shifting the rotation parameters:
\[
\theta_a \mapsto \theta_a + q A_a,
\]
yielding the extended rotor
\[
R_{\text{dual}}^{\text{EM}} = \exp\left( \frac{1}{2} \sum_a (\theta_a + q A_a) \Gamma_a + \frac{q}{2} \int F_{\text{dual}} \, ds \right),
\]
where $F_{\text{dual}} = \sum_a (E_a i \Gamma_a + B_a \Gamma_a)$ is the electromagnetic bivector. The Lorentz force emerges from the sandwich product $\tilde{X}' = R_{\text{dual}}^{\text{EM}} X' (R_{\text{dual}}^{\text{EM}})^\dagger \approx X' + q F_{\text{dual}} \wedge X'$.

The weak force is incorporated via chiral asymmetry, coupling only to left-handed components:
\[
\theta_a^L = \theta_a + g_W W_a, \quad \theta_a^R = \theta_a,
\]
with the chiral rotor
\[
R_{\text{dual}}^{\text{weak}} = \exp\left( \frac{1}{2} \sum_a \theta_a^L \Gamma_a P_L + \frac{1}{2} \sum_a \theta_a^R \Gamma_a P_R \right),
\]
where $W_a$ denotes weak gauge bivectors. The V-A structure follows from the time-reversed dual sector, reproducing parity violation.

Unification occurs in the full torsional bivector:
\[
\Omega_{\text{biv}} = \log\left( R_{\text{usual}}^\dagger R_{\text{dual}}^{\text{EM+weak}} \right) \approx \sum_a (\theta_a + q A_a + g_W W_a P_L - \phi_a i) \Gamma_a.
\]
The scalar $J = \frac{1}{16} B(\Omega_{\text{biv}}, \Omega_{\text{biv}})$ now includes electromagnetic and weak contributions, with the electroweak symmetry emerging as a subgroup of $\mathrm{Spin}^+(3,1) \oplus \mathrm{Spin}^+(3,1)$. Mass generation via rotor rigidity unifies with the Higgs mechanism, eliminating fine-tuning.

This embedding recovers the Standard Model electroweak sector at low energies while resolving hierarchy problems through dual compactness, paving the way for full unification with strong and gravitational forces.

\section{Gravitational Engineering: Coherent Excitation of Dual-Spacetime Rotors}
\label{sec:gravitycontrol}

The double spacetime theory transforms gravity from an inviolable geometric constraint into an engineerable torsional degree of freedom. The torsion scalar $J = \frac{1}{16} B(\Omega_{\text{biv}}, \Omega_{\text{biv}})$, with $\Omega_{\text{biv}} = \log(R_{\text{usual}}^\dagger R_{\text{dual}})$, depends solely on the relative misalignment between the particle-intrinsic usual and dual rotors. External fields capable of coherently modulating the dual rotor parameters $\theta_a$ can therefore reduce $J \to 0$ (gravitational shielding or inertial mass reduction) or drive $J < 0$ (repulsive gravity).

\subsection{Theoretical Coupling of Circularly Polarized Electromagnetic Fields to Dual Rotors}

The electromagnetic field is naturally partitioned across the two spacetimes:

\[
F = F_{\text{usual}} + F_{\text{dual}},
\]

where

\[
F_{\text{usual}} = \mathbf{E} \cdot (iI, iJ, iK) \quad (\hat{p}^2 = +1:\ \text{boost-like, non-compact}),
\]

\[
F_{\text{dual}} = \mathbf{B} \cdot (I, J, K) \quad (\hat{q}^2 = -1:\ \text{rotation-like, compact}).
\]

The electric component $\mathbf{E}$ couples primarily to the usual-spacetime boost generators, while the magnetic component $\mathbf{B}$ couples exclusively to the dual-spacetime rotation generators. This algebraic separation is the geometric origin of the distinct transformation properties under time reversal ($T$: $\mathbf{E} \to \mathbf{E}$, $\mathbf{B} \to -\mathbf{B}$): the dual spacetime is intrinsically time-reversed.

For a circularly polarized wave propagating along $z$ with helicity $\sigma = \pm 1$,

\[
\mathbf{E} = E_0 (\hat{x} \cos(kz - \omega t) + \sigma \hat{y} \sin(kz - \omega t)), \quad \mathbf{B} = \frac{E_0}{c} (-\hat{x} \sin(kz - \omega t) + \sigma \hat{y} \cos(kz - \omega t)).
\]

Time-averaging eliminates the oscillatory boost-like contribution of $\mathbf{E}$, leaving a pure dual-spacetime rotational excitation:

\[
\langle F_{\text{dual}} \rangle \propto \sigma E_0^2 \, K.
\]

The helicity $\sigma$ directly determines the direction of dual-rotor drive: right-handed polarization ($\sigma = +1$) reinforces attractive torsional mismatch ($J > 0$), while left-handed polarization ($\sigma = -1$) opposes it, reducing or reversing the scalar $J$. This helicity selectivity is the sharpest manifestation of the theory: gravity is not merely modulated but chirally addressed via the intrinsic time-reversal asymmetry of the dual spacetime.

The coupling is minimal and gauge-invariant:

\[
R_{\text{dual}} \to \mathcal{P} \exp\!\left( \frac{1}{2} \int (\theta_a \Gamma_a + e \lambda A_a \Gamma_a) \, ds \right),
\]

where the dual generators $\Gamma_a$ mediate the interaction. Near resonance, the dual phase evolves as

\[
\dot{\theta} = \kappa \sigma E_0^2 \sin(\omega t - \theta + \delta_\sigma),
\]

yielding coherent excitation of the torsional mismatch angle.

\subsection{Resonance Frequencies and the Role of Superconductors}

The intrinsic dual-spacetime scale is Planckian ($\sim 10^{43}$ Hz). Macroscopic coherence ($N \sim 10^{26}$ particles) redshifts the effective frequency to $10^{12}$--$10^{15}$ Hz. Superconductors further reduce it dramatically via Josephson plasma resonances and collective phase locking, achieving tunable modes in the GHz--THz range accessible to current technology.

\subsection{Proposed Experimental Protocol and Predictions}

A microgram-scale test mass in a high-$Q$ superconducting resonator is illuminated with tunable circularly polarized radiation (0.1--10 THz). Expected signatures include resonant, helicity-dependent changes in apparent weight or inertial mass ($\Delta m/m \gtrsim 10^{-6}$), with left-handed polarization producing reduction and right-handed enhancement. Success at this level scales directly to practical gravitational engineering.

For the first time, gravity is rendered not merely controllable but chirally controllable, with helicity serving as the switch between attraction and repulsion in the torsional geometry of double spacetime.

\section{Baryon Asymmetry: The Geometric Origin from Dual Rotor Time Arrow Fixation}
\label{sec:baryonasymmetry}

In the double spacetime theory, the observed baryon asymmetry of the universe---the profound imbalance where baryons vastly outnumber antibaryons, quantified by the parameter $\eta \equiv (n_B - n_{\bar{B}})/n_\gamma \approx 6 \times 10^{-10}$---emerges as an inevitable consequence of the asymmetric fixation of the time arrow in the dual rotors during the Planck epoch. This mechanism, rooted in the intrinsic chirality of the biquaternion algebra $\mathrm{Cl}(3,1)$, satisfies all Sakharov conditions without invoking additional fields, parameters, or fine-tuning. Here, we derive the asymmetry algebraically, emphasizing the geometric interplay between the usual and dual spacetimes.

\subsection{The Intrinsic Time Duality and Rotor Parameters}

Recall that each particle encodes a pair of Minkowski spacetimes within the 16-real-dimensional biquaternionic structure: the usual spacetime spanned by $\{j, kI, kJ, kK\}$ with future-directed time basis $j$, and its dual counterpart $\{k, jI, jJ, jK\}$ obtained via right-multiplication by the pseudoscalar $i$, yielding $j i = -k$ and thus a past-directed time basis $k$. This duality operation $X \mapsto X i$ preserves the Minkowski norm (see Sec.~\ref{sec:math}, where the explicit non-commutative computation gives $(X i)^2 = X^2$) but inverts the temporal orientation, imprinting a fundamental chirality: particles (usual-dominant) propagate forward in time, while antiparticles (dual-dominant) would inherently propagate backward if not for dynamical suppression.

The full kinematic evolution is governed by the complete rotor. In vacuum, torsionlessness ($J=0$) is required, ensuring perfect anti-synchronization of the time arrows: the forward usual rotor is mirrored by a retrograde dual rotor, maintaining particle-antiparticle equivalence.

\subsection{Planck-Epoch Fixation of the Time Arrow}

At the Planck scale ($t \sim 10^{-43}$ s, $T \sim 10^{19}$ GeV), quantum gravitational fluctuations---manifest as zero-point oscillations in the dual rotor bivectors---induce a spontaneous symmetry breaking via a phase transition in the vacuum expectation value of the rotor parameters. The biquaternionic basis asymmetry, where $j$ is the sole time-like generator while spatial bases involve cross-terms ($kI$, etc.), selects a unique energetically favorable configuration:

\[
\langle \theta \rangle = \phi + \delta, \quad \delta > 0,
\]
where $\delta$ is the CP-odd chiral phase imprinted at the Planck epoch. We take $\delta$ as a small, free parameter of the framework whose magnitude is to be fixed by observation; the rough estimate $\delta \sim 10^{-5}$ (comparable to the Jarlskog measure of the observed CKM phase) is adopted throughout. This fixation arises because the pseudoscalar $i$ commutes with $\Gamma_a$ but anticommutes with $i \Gamma_a$, generating a chiral potential
\[
V_\text{chiral} = \frac{1}{2} \mathrm{Tr} \left[ (i \Gamma_a) [\Gamma_b, i] \right] \propto \epsilon_{abc} \delta_a \delta_b \delta_c,
\]
which minimizes when the dual rotor ``over-rotates'' in the positive direction relative to the usual rotor. Consequently, the torsional mismatch bivector becomes
\[
\Omega_\text{biv} = \log(R_\text{usual}^\dagger R_\text{dual}) \approx \sum_a \delta_a (i \Gamma_a + \Gamma_a) \in \mathfrak{so}(3,1) \oplus \mathfrak{so}(3,1),
\]
yielding a positive torsion scalar $J = \frac{1}{16} B(\Omega_\text{biv}, \Omega_\text{biv}) > 0$. This breaks the particle-antiparticle symmetry at the algebraic core: usual-dominant modes ($\phi$) experience attractive torsional binding ($J > 0$), while dual-dominant modes ($\theta$) are repelled into $J < 0$ layers.

\subsection{Sakharov Conditions and the Order-of-Magnitude Estimate}

The fixation automatically fulfills the three Sakharov conditions through the rotor exponential:

1. Baryon Number Violation: Baryon number $B$ is encoded as the pseudoscalar phase in the dual rotor, $B \propto \arg(\det R_\text{dual})$. The mismatch $\delta \neq 0$ enables non-conserving processes such as $X \mapsto X i i = -X$ (double time reversal, effectively a baryon flip), so the rates for usual- and dual-dominant modes differ at $O(\delta)$.

2. C and CP Violation: The dual map $i$ is intrinsically chiral (odd under parity, as $i \mapsto -i$ under spatial inversion), and $\delta > 0$ introduces a CP-odd phase $\phi_\text{CP} = \arg(\exp(i \delta \Gamma_a)) \sim \delta$, comparable to the observed CKM phase. This geometric CP violation is universal, permeating all fermion generations without Yukawa-specific tuning.

3. Departure from Thermal Equilibrium: Inflationary expansion drives the system far from equilibrium, allowing the CP-odd phase $\delta$ to be converted into a frozen-in number-density asymmetry as dual-dominant modes annihilate against their usual-dominant partners. A standard freeze-out estimate gives an asymmetry that is quadratic in the chiral phase,
\[
\eta \sim \delta^2,
\]
so that the observed value $\eta \sim 6 \times 10^{-10}$ is reproduced by $\delta \sim 10^{-5}$, in agreement with the CP-odd phase estimate above. The precise numerical coefficient depends on the post-inflationary reheating temperature and is not predicted at this level.

Post-inflation, residual dual-dominant antiparticles annihilate against the baryon surplus, leaving $\eta$ as the frozen relic of the time arrow fixation. Antiparticles do not ``disappear'' but are dynamically excluded: their backward time arrow is incompatible with the forward-directed cosmic expansion fixed by $\delta > 0$.

\subsection{Observational Signatures and Falsifiability}

This mechanism predicts subtle CMB anisotropies from chiral gravitational waves sourced by $\Omega_\text{biv}$, with a tensor-to-scalar ratio $r \sim \delta^2 \sim 10^{-10}$ and a CP-odd parity spectrum detectable by future missions like LiteBIRD. At low energies, it implies a universal lepton asymmetry $\eta_\ell \approx \eta_B$, resolvable by neutrinoless double-beta decay experiments (e.g., LEGEND-200). Non-observation of these signatures would falsify the theory, while confirmation would elevate dual spacetime to the foundational framework for cosmology.

In essence, the baryon asymmetry is not a puzzle but the geometric echo of the universe's choice of time's arrow: forward for matter, forbidden for antimatter. The dual rotors, once fixed, decree that our cosmos is a sanctuary for particles alone.

\section{Discrete Time and Compactified Torsion Scalar}
\label{sec:discretetime}

The rejection of the spacetime continuum hypothesis demands a discrete structure at the fundamental scale. The intrinsic duality encoded in the biquaternion algebra—particularly the pseudoscalar $i$ and the dual map $X \mapsto Xi$ that reverses the temporal arrow while preserving the Minkowski norm—compactifies the dual spacetime intrinsically, while leaving the usual spacetime non-compact.

We therefore discretize only the dual-spacetime rotation angle:
\[
\theta \in \frac{2\pi}{N} \mathbb{Z},
\]
where $N \gg 1$ is a very large integer determining the compactification scale of the dual spacetime (typically proportional to the particle’s rest mass in Planck units). The usual-spacetime rapidity remains continuous:
\[
\phi \in \mathbb{R}.
\]

This asymmetry is required for theoretical consistency:
\begin{itemize}
  \item The dual spacetime is elliptic ($\hat{q}^2 = -1$), inherently compact, and its discretization directly reflects the $4\pi$-periodic topology of $S^3$ parameterized by $R_{\text{dual}}$. This provides the finite phase volume essential for resolving the Newton–Wigner localization problem and yielding natural quantum discreteness.
  \item The usual spacetime is hyperbolic ($\hat{p}^2 = +1$), intrinsically non-compact, so continuity of $\phi$ preserves the full Lorentz group and ensures exact equivalence with classical general relativity in the macroscopic limit.
\end{itemize}

The torsional mismatch rotor $\Omega = R_{\text{usual}}^\dagger R_{\text{dual}}$ thus combines a continuous usual component with a discrete dual component. The torsion scalar $J = \frac{1}{16} B(\Omega_{\text{biv}}, \Omega_{\text{biv}})$ remains bounded, rigorously excluding continuum-induced singularities.

A critical issue arises in the effective Lorentz factor across torsional layers, which typically appears as a synchronization function of the form
\[
\gamma_{\text{eff}} \propto D_N(\delta\theta) \equiv \frac{\sin(N \delta\theta / 2)}{N \sin(\delta\theta / 2)},
\]
where $\delta\theta$ is the torsional mismatch angle and $D_N$ is the (normalized) Dirichlet kernel. The denominator vanishes precisely on the lattice $\delta\theta \in 2\pi \mathbb{Z}$.

At every such point the numerator $\sin(N \cdot k\pi)$ also vanishes (for any integer $N$, prime or composite), so the apparent singularity is in fact a removable $0/0$ of Dirichlet-kernel type. The standard L'H\^{o}pital evaluation gives the finite limit
\[
\lim_{\delta\theta \to 2\pi k} D_N(\delta\theta) = \pm 1 \qquad (k \in \mathbb{Z}),
\]
so the synchronization function is bounded on the entire compact lattice and the classical continuum limit $N \to \infty$ is recovered without pathology.

Thus the boundedness of the torsional dynamics is a structural property of the Dirichlet kernel and does not require $N$ to be prime; any large positive integer $N$ suffices to suppress spurious singularities while preserving the continuum limit. The number-theoretic structure of the dual-spacetime compactification---in particular, the relation between $N$ and particle mass spectra---remains a natural extension of the framework and is developed in detail in a companion paper on the discrete dual spacetime model.

\section{Dual-Rotor Dynamics and the Geometric Origin of the de Broglie Phase}
\label{sec:dualrotordynamics}

In the double spacetime theory, the torsional mismatch scalar $J$ is constructed from the relative rotor $\Omega = R_{\text{usual}}^\dagger R_{\text{dual}}$. While $J$ serves as the gravitational action density in the classical field theory, it also admits a natural particle-intrinsic interpretation as a potential governing the internal degrees of freedom of individual massive particles.

For small mismatch angles, the Killing form invariant expands as
\[
J \approx \frac{1}{2} \sum_{a=1}^3 \delta\phi_a^2,
\]
where $\delta\phi_a = \phi_a - \theta_a$ is the torsional mismatch angle in each boost-rotation plane. Incorporating the rest mass $m$ as the rigidity constant of the dual sector yields
\[
J = \frac{1}{2} m \sum_{a=1}^3 \delta\phi_a^2.
\]

We propose an effective action for the intrinsic rotor angles of a single particle:
\[
S_{\text{particle}} = \int \left[ \frac{1}{2} \sum_{a=1}^3 \dot{\phi}_a^2 - \frac{1}{2} \sum_{a=1}^3 \dot{\theta}_a^2 - \frac{m}{2} \sum_{a=1}^3 (\phi_a - \theta_a)^2 \right] dt.
\]
The opposite signs of the kinetic terms reflect the usual sector's boost-like character and the dual sector's time-reversed rotation-like character.

The Euler--Lagrange equations are
\begin{align}
\ddot{\phi}_a &= m (\phi_a - \theta_a), \\
-\ddot{\theta}_a &= m (\phi_a - \theta_a).
\end{align}
Defining the torsional mismatch angle $\delta\phi_a = \phi_a - \theta_a$, we obtain the simple harmonic oscillator equation
\begin{equation}
\ddot{\delta\phi}_a + m \delta\phi_a = 0.
\label{eq:dual-harmonic}
\end{equation}
The oscillation frequency is $\sqrt{m}$ (in natural units where $\hbar = c = 1$).

For a free particle in uniform motion, the usual rapidity components evolve with constant rate $\dot{\phi}_a = \gamma v_a/c$. The general solution of Eq.~\eqref{eq:dual-harmonic} is
\[
\delta\phi_a(t) = A_a \sin(\sqrt{m}\, t + \psi_a).
\]
In the low-energy, non-relativistic regime, rapid Compton-frequency oscillations average to zero, leaving a slow linear accumulation of phase proportional to the proper time. Restoring $\hbar$ and $c$, the accumulated phase becomes
\[
\int \delta\phi_a \, dt \propto \frac{E t - \mathbf{p} \cdot \mathbf{x}}{\hbar},
\]
which is precisely the de Broglie phase of the matter wave.

Thus, the quantum mechanical wave function phase $S(x,t)/\hbar$ emerges geometrically as the slowly varying component of the dual-rotor lag relative to the usual rotor. This provides a fully background-independent, algebraic origin for quantum matter waves within the double spacetime framework, without invoking additional postulates.

Collective averaging of this particle-intrinsic action over many particles recovers the classical teleparallel action $\int J \, d^4x$, ensuring complete consistency with general relativity at macroscopic scales while revealing the microscopic quantum dynamics of the dual spacetime degrees of freedom.

This harmonic oscillator dynamics constitutes the bridge between the classical torsional description of gravity and the full quantum theory developed in the following section.

\section{Towards Quantum Gravity and Full Unification}
\label{sec:quantumgravity}

The classical formulation presented thus far demonstrates that gravity emerges as torsional mismatch between the particle-intrinsic usual and dual spacetimes, encoded in the 16-dimensional biquaternion algebra isomorphic to $\mathrm{Cl}(3,1)$. The action $S = \frac{c^4}{16\pi G} \int J \, d^4x$, with
\[
J = \frac{1}{16} B(\Omega_\text{biv}, \Omega_\text{biv}), \quad \Omega_\text{biv} = \log(R_\text{usual}^\dagger R_\text{dual}),
\]
yields exact equivalence to the Einstein-Hilbert action while abandoning the continuum hypothesis and Riemannian curvature.

A profound extension arises from recognizing that the dual rotor $R_\text{dual}$ is mathematically identical to the SU(2) spin-1/2 rotation operator. The dual-sector phases $\theta$ thus directly parameterize quantum states. Due to the intrinsic compactness of the dual spacetime ($\theta \mod 4\pi$), these phases live on a compact manifold, naturally yielding a Hilbert space structure where each particle's quantum state $|\psi\rangle$ corresponds to the unitary phase part of $R_\text{dual}$.

Position and momentum operators emerge as projections of dual rotor angles, while forces—including gravity itself—manifest as modulations of relative rotor phases. This identifies quantum mechanics with dual spacetime rotations and gravity with torsional mismatch in the same algebraic framework, eliminating singularities via bounded discrete torsion layers.

The full quantum gravity framework, including explicit Hilbert space construction, spinor embedding in minimal left ideals of $\mathrm{Cl}(3,1)$, unification of all interactions, and resolution of baryon asymmetry through Planck-era time-arrow fixation, is developed in detail in the companion paper ``Quantum Gravity as Torsional Mismatch in Biquaternion Double Spacetime'' \cite{dst-quantum-gravity}.

This extension completes the theory: General Relativity is preserved classically, yet quantum gravity emerges algebraically without extra dimensions or ad-hoc quantization—solely from the double spacetime structure intrinsic to every massive particle.

\section{Conclusions}
\label{sec:conclusions}
The double spacetime theory, grounded in the biquaternionic algebra $\mathbb{H} \otimes_{\mathbb{R}} \mathbb{H} \cong \mathrm{Cl}(3,1)$, offers a transformative framework for understanding gravity, inertia, and fundamental interactions. By embedding each particle within a pair of intertwined Minkowski spacetimes—one usual and one dual—the theory reinterprets gravitational phenomena as manifestations of torsional misalignment between these sectors. This paradigm shift resolves longstanding paradoxes in gravitational physics, unifies the strong nuclear force as a manifestation of layered torsional dynamics, and provides a natural mechanism for baryon asymmetry through intrinsic chirality and time arrow fixation. The theory's predictions, from gravity control via dual rotor synchronization to the reinterpretation of atomic nuclei as primordial torsion stars, open new avenues for experimental verification and technological innovation. As we stand on the cusp of a new era in gravitational physics, the dual spacetime framework beckons us to rethink our understanding of the cosmos and our place within it.

\bibliographystyle{unsrt}
\begin{thebibliography}{1}

\bibitem{einstein1928}
A.~Einstein,
``Riemann-Geometrie mit Aufrechterhaltung des Begriffes des Fernparallelismus,''
\textit{Sitzungsber. Preuss. Akad. Wiss. Phys.-Math. Kl.},
217--221 (1928).

\bibitem{einstein1939}
A.~Einstein,
``On a stationary system with spherical symmetry consisting of many gravitating masses,''
\textit{Ann. Math.},
\textbf{40}, no.~4, 922--936 (1939).
doi:10.2307/1968902

\bibitem{girard1999}
P.~R. Girard,
``Einstein’s equations and Clifford algebra,''
\textit{Advances in Applied Clifford Algebras},
\textbf{9}, no.~2, 225--230 (1999).
doi:10.1007/BF03042377

\bibitem{maluf2013}
J.~W. Maluf,
``The teleparallel equivalent of general relativity,''
\textit{Ann. Phys. (Berlin)},
\textbf{525}, 339--359 (2013).
doi:10.1002/andp.201200272

\bibitem{dst-quantum-gravity} Anonymous 2025, ``Quantum Gravity as Torsional Mismatch in Biquaternion Double Spacetime'', https://github.com/hypernumbernet/dual-spacetime-doc/blob/main/dst-quantum-gravity.pdf.

\end{thebibliography}

\clearpage
\appendix
\section{Multiplication Table}
\label{app:mult-table}

\begin{table}[ht]
  \begin{center}
    \caption{Biquaternion Multiplication Table}
    \small
    \begin{tabular}{r|rrrrrrrrrrrrrrr}
      &$j$&$kI$&$kJ$&$kK$&$iI$&$iJ$&$iK$&$I$&$J$&$K$&$k$&$jI$&$jJ$&$jK$&$i$ \\ \hline
      j&$-1$&$+iI$&$+iJ$&$+iK$&$-kI$&$-kJ$&$-kK$&$+jI$&$+jJ$&$+jK$&$+i$&$-I$&$-J$&$-K$&$-k$ \\ \hline
      kI&$-iI$&$+1$&$-K$&$+J$&$-j$&$+jK$&$-jJ$&$-k$&$+kK$&$-kJ$&$-I$&$+i$&$-iK$&$+iJ$&$+jI$ \\ \hline
      kJ&$-iJ$&$+K$&$+1$&$-I$&$-jK$&$-j$&$+jI$&$-kK$&$-k$&$+kI$&$-J$&$+iK$&$+i$&$-iI$&$+jJ$ \\ \hline
      kK&$-iK$&$-J$&$+I$&$+1$&$+jJ$&$-jI$&$-j$&$+kJ$&$-kI$&$-k$&$-K$&$-iJ$&$+iI$&$+i$&$+jK$ \\ \hline
      iI&$+kI$&$+j$&$-jK$&$+jJ$&$+1$&$-K$&$+J$&$-i$&$+iK$&$-iJ$&$-jI$&$-k$&$+kK$&$-kJ$&$-I$ \\ \hline
      iJ&$+kJ$&$+jK$&$+j$&$-jI$&$+K$&$+1$&$-I$&$-iK$&$-i$&$+iI$&$-jJ$&$-kK$&$-k$&$+kI$&$-J$ \\ \hline
      iK&$+kK$&$-jJ$&$+jI$&$+j$&$-J$&$+I$&$+1$&$+iJ$&$-iI$&$-i$&$-jK$&$+kJ$&$-kI$&$-k$&$-K$ \\ \hline
      I&$+jI$&$-k$&$+kK$&$-kJ$&$-i$&$+iK$&$-iJ$&$-1$&$+K$&$-J$&$+kI$&$-j$&$+jK$&$-jJ$&$+iI$ \\ \hline
      J&$+jJ$&$-kK$&$-k$&$+kI$&$-iK$&$-i$&$+iI$&$-K$&$-1$&$+I$&$+kJ$&$-jK$&$-j$&$+jI$&$+iJ$ \\ \hline
      K&$+jK$&$+kJ$&$-kI$&$-k$&$+iJ$&$-iI$&$-i$&$+J$&$-I$&$-1$&$+kK$&$+jJ$&$-jI$&$-j$&$+iK$ \\ \hline
      k&$-i$&$-I$&$-J$&$-K$&$+jI$&$+jJ$&$+jK$&$+kI$&$+kJ$&$+kK$&$-1$&$-iI$&$-iJ$&$-iK$&$+j$ \\ \hline
      jI&$-I$&$-i$&$+iK$&$-iJ$&$+k$&$-kK$&$+kJ$&$-j$&$+jK$&$-jJ$&$+iI$&$+1$&$-K$&$+J$&$-kI$ \\ \hline
      jJ&$-J$&$-iK$&$-i$&$+iI$&$+kK$&$+k$&$-kI$&$-jK$&$-j$&$+jI$&$+iJ$&$+K$&$+1$&$-I$&$-kJ$ \\ \hline
      jK&$-K$&$+iJ$&$-iI$&$-i$&$-kJ$&$+kI$&$+k$&$+jJ$&$-jI$&$-j$&$+iK$&$-J$&$+I$&$+1$&$-kK$ \\ \hline
      i&$+k$&$-jI$&$-jJ$&$-jK$&$-I$&$-J$&$-K$&$+iI$&$+iJ$&$+iK$&$-j$&$+kI$&$+kJ$&$+kK$&$-1$ \\ \hline
    \end{tabular}
  \end{center}
\end{table}
\begin{table}[ht]
  \begin{center}
    \caption{Cl(3,1) Variant Multiplication Table ($e_0, e_1, e_2, e_3 \mapsto 0,1,2,3$)}
    \scriptsize
    \begin{tabular}{r|rrrrrrrrrrrrrrr}
      &$0$&$1$&$2$&$3$&$01$&$02$&$03$&$32$&$13$&$21$&$123$&$032$&$013$&$021$&$0123$\\\hline
      $0$&$-$&$+01$&$+02$&$+03$&$-1$&$-2$&$-3$&$+032$&$+013$&$+021$&$+0123$&$-32$&$-13$&$-21$&$-123$\\
      $1$&$-01$&$+$&$-21$&$+13$&$-0$&$+021$&$-013$&$-123$&$+3$&$-2$&$-32$&$+0123$&$-03$&$+02$&$+032$\\
      $2$&$-02$&$+21$&$+$&$-32$&$-021$&$-0$&$+032$&$-3$&$-123$&$+1$&$-13$&$+03$&$+0123$&$-01$&$+013$\\
      $3$&$-03$&$-13$&$+32$&$+$&$+013$&$-032$&$-0$&$+2$&$-1$&$-123$&$-21$&$-02$&$+01$&$+0123$&$+021$\\
      $01$&$+1$&$+0$&$-021$&$+013$&$+$&$-21$&$+13$&$-0123$&$+03$&$-02$&$-032$&$-123$&$+3$&$-2$&$-32$\\
      $02$&$+2$&$+021$&$+0$&$-032$&$+21$&$+$&$-32$&$-03$&$-0123$&$+01$&$-013$&$-3$&$-123$&$+1$&$-13$\\
      $03$&$+3$&$-013$&$+032$&$+0$&$-13$&$+32$&$+$&$+02$&$-01$&$-0123$&$-021$&$+2$&$-1$&$-123$&$-21$\\
      $32$&$+032$&$-123$&$+3$&$-2$&$-0123$&$+03$&$-02$&$-$&$+21$&$-13$&$+1$&$-0$&$+021$&$-013$&$+01$\\
      $13$&$+013$&$-3$&$-123$&$+1$&$-03$&$-0123$&$+01$&$-21$&$-$&$+32$&$+2$&$-021$&$-0$&$+032$&$+02$\\
      $21$&$+021$&$+2$&$-1$&$-123$&$+02$&$-01$&$-0123$&$+13$&$-32$&$-$&$+3$&$+013$&$-032$&$-0$&$+03$\\
      $123$&$-0123$&$-32$&$-13$&$-21$&$+032$&$+013$&$+021$&$+1$&$+2$&$+3$&$-$&$-01$&$-02$&$-03$&$+0$\\
      $032$&$-32$&$-0123$&$+03$&$-02$&$+123$&$-3$&$+2$&$-0$&$+021$&$-013$&$+01$&$+$&$-21$&$+13$&$-1$\\
      $013$&$-13$&$-03$&$-0123$&$+01$&$+3$&$+123$&$-1$&$-021$&$-0$&$+032$&$+02$&$+21$&$+$&$-32$&$-2$\\
      $021$&$-21$&$+02$&$-01$&$-0123$&$-2$&$+1$&$+123$&$+013$&$-032$&$-0$&$+03$&$-13$&$+32$&$+$&$-3$\\
      $0123$&$+123$&$-032$&$-013$&$-021$&$-32$&$-13$&$-21$&$+01$&$+02$&$+03$&$-0$&$+1$&$+2$&$+3$&$-$\\
    \end{tabular}
  \end{center}
\end{table}
These are indeed the same results and are isomorphic.
\clearpage
\section{Equivalence of the Torsional Mismatch Scalar J to the Teleparallel Torsion Scalar T of TEGR}
\label{app:TEGR-equivalence}

The torsional mismatch bivector is
\[
\Omega_{\rm biv} = \log(R_{\rm usual}^\dagger R_{\rm dual}) = \sum_{a=1}^3 \left( \frac{\alpha_a}{2} i\Gamma_a + \frac{\beta_a}{2} \Gamma_a \right),
\]
with \( (i\Gamma_a)^2 = +1 \), \( \Gamma_a^2 = -1 \), and \([i\Gamma_a, \Gamma_b] = 0\).

The Killing form on \(\mathfrak{so}(3,1) \oplus \mathfrak{so}(3,1)\) satisfies \( B(i\Gamma_a, i\Gamma_a) = +8 \), \( B(\Gamma_a, \Gamma_a) = -8 \), yielding
\[
B(\Omega_{\rm biv}, \Omega_{\rm biv}) = 8 \sum_{a=1}^3 (\alpha_a^2 - \beta_a^2).
\]
Thus the torsional mismatch scalar is
\[
J = \frac{1}{16} B(\Omega_{\rm biv}, \Omega_{\rm biv}) = \frac12 \sum_{a=1}^3 (\alpha_a^2 - \beta_a^2).
\]

To compare this invariant with TEGR we must first specify the geometric correspondence between the rotor field and the tetrad. Concretely, the total rotor field $R_{\rm total}(x) = R_{\rm usual}(x)\, R_{\rm dual}(x)$ acts on a fixed reference tetrad $\{\mathring{e}^a\}$ by the spinor sandwich product
\[
e^a{}_\mu(x) \;=\; \langle R_{\rm total}(x)\, \mathring{e}^a\, R_{\rm total}^\dagger(x) \rangle_{1,\mu},
\]
where $\langle \cdot \rangle_{1,\mu}$ extracts the grade-1 (vector) component along the coordinate basis. This identification promotes $R_{\rm total}$ to a tetrad field $e^a{}_\mu$ with metric $g_{\mu\nu} = \eta_{ab}\, e^a{}_\mu e^b{}_\nu$, and renders the dual-mismatch bivector $\Omega_{\rm biv}$ proportional to the contorsion of the Weitzenb\"ock connection associated with $e^a{}_\mu$:
\[
\Omega_{\rm biv}\big|_{\mu\nu}^{ab} \;\propto\; K^{ab}{}_\mu\, dx^\mu \wedge dx^\nu,
\qquad
T^a{}_{\mu\nu} \;=\; \partial_\mu e^a{}_\nu - \partial_\nu e^a{}_\mu.
\]
Under this identification the quadratic invariant $J$ corresponds (up to a total divergence) to the teleparallel torsion scalar
\[
T = \frac14 T_{\rho\mu\nu}T^{\rho\mu\nu} + \frac12 T_{\rho\mu\nu}T^{\nu\mu\rho} - T_\rho T^\rho,
\]
giving the on-shell identification
\[
J = \frac12 T + \nabla_\mu B^\mu,
\]
where $\nabla_\mu B^\mu$ is a total divergence whose explicit form is fixed by the projection above.

It is a standard result of TEGR \cite{maluf2013} that, for the same underlying tetrad,
\[
\int T\, e\, d^4x = -\int R\sqrt{-g}\, d^4x + \int \partial_\mu(\sqrt{-g}\, V^\mu)\, d^4x,
\]
i.e.\ $T$ and $-R$ differ only by a total divergence. Substituting the relation for $J$ gives
\[
S = \frac{c^4}{16\pi G} \int J \, d^4x = -\frac{c^4}{16\pi G} \int R \sqrt{-g}\, d^4x + (\text{boundary terms}).
\]

The variation of $S$ is taken with the tetrad $e^a{}_\mu$ (or equivalently the rotor field $R_{\rm total}$) as the independent dynamical variable, with the standard boundary condition that all field variations $\delta e^a{}_\mu$ vanish on the boundary $\partial\mathcal{M}$ of the integration region. Under this condition the total-divergence terms above do not contribute to the variational equations, and the variation of the remainder reproduces the Einstein equations
\[
G_{\mu\nu} = \frac{8\pi G}{c^4} T_{\mu\nu}.
\]

Thus, with the rotor-to-tetrad identification stated above and the standard variational boundary condition, the Double Spacetime action is dynamically equivalent to the Einstein--Hilbert action in the classical limit.

\end{document}
